% compile in holomorphic-twist-program/notes/flagship:  pdflatex flagship.tex  (twice)
\documentclass[10pt,a4paper]{article}
\usepackage[T1]{fontenc}
\usepackage[utf8]{inputenc}
\usepackage{lmodern}
\usepackage[margin=2.1cm]{geometry}
\usepackage{amsmath,amssymb,amsthm,mathtools}
\usepackage{booktabs,longtable,array}
\usepackage[colorlinks=true,linkcolor=blue!50!black,citecolor=blue!50!black,urlcolor=blue!50!black]{hyperref}
\usepackage{xcolor}
\numberwithin{equation}{section}
\allowdisplaybreaks
\setlength{\parindent}{0pt}
\setlength{\parskip}{3pt}
\emergencystretch=3em

\newtheorem{theorem}{Theorem}[section]
\newtheorem{proposition}[theorem]{Proposition}
\newtheorem{lemma}[theorem]{Lemma}
\newtheorem{corollary}[theorem]{Corollary}
\theoremstyle{definition}
\newtheorem{definition}[theorem]{Definition}
\newtheorem{dstatement}[theorem]{Derived statement}
\theoremstyle{remark}
\newtheorem{remark}[theorem]{Remark}

\newcommand{\RPp}[1]{\mathbb{RP}^{#1}_{>0}}
\newcommand{\Pf}{\operatorname{Pf}}
\newcommand{\tr}{\operatorname{tr}}
\newcommand{\Tr}{\operatorname{Tr}}
\newcommand{\alt}{\operatorname{alt}}
\newcommand{\sgn}{\operatorname{sgn}}
\newcommand{\g}{\mathfrak g}
\newcommand{\fsl}{\mathfrak{sl}}
\newcommand{\fgl}{\mathfrak{gl}}
\newcommand{\fso}{\mathfrak{so}}
\newcommand{\fsp}{\mathfrak{sp}}
\newcommand{\dbar}{\bar\partial}
\newcommand{\zb}{\bar z}
\newcommand{\cN}{\mathcal N}
\newcommand{\Aut}{\operatorname{Aut}}
\newcommand{\htk}{\mathrm{htk}}
\newcommand{\Vht}{V_{\rm ht}}
\newcommand{\Veta}{V^{\eta}}
\newcommand{\Cg}{C^0_{\g}}
\newcommand{\cert}[1]{\par\noindent\textup{Certificates: #1.}\par}
\newcommand{\uf}{u_2f}
\newcommand{\Dthree}{D_3}

\title{Commutative Chern--Simons theory on $\mathbb R\times\mathbb C^2$ through order $\hbar^3$\\ and the two-loop theta kernels on $\mathbb R^T\times\mathbb C^H$}
\author{}
\date{2026-11-21}

\begin{document}
\maketitle

\begin{abstract}
\noindent Two-loop theta kernels on $\mathbb R^T\times\mathbb C^H$ ($T\le1$, $H+T\le6$) in closed form. Commutative Chern--Simons theory on $\mathbb R\times\mathbb C^2$, gauge algebra $\g$ with invariant form $\kappa$, heat-kernel gauge. (1) Translation-invariant local functionals annihilated by $\iota_{\partial_\tau},\iota_{\partial_{\zb_1}},\iota_{\partial_{\zb_2}}$ have the cohomology of Chevalley--Eilenberg cochains of $\g[[z_1,z_2]]$ on $\sum p=0$, ghost number $g\leftrightarrow$ degree $g+3$; without these conditions there are $8$ summands, and the heat-kernel scheme satisfies them. (2) The two-loop obstruction is $O_2=\frac{\pi\cN}{240}\int\Omega[\partial_1\partial_2][\partial_3\partial_4](-3W_{133}+2W_{223})$, $\cN=-U\hbar^2/(64\pi^5)$, $U=\pm1$; it is not exact for every simple $\g\neq\fsl_2$ and for $\fgl_N$ ($\kappa=\tr$, $N\ge3$), and exact for $\fsl_2,\fgl_1,\fgl_2$. (3) $U(2)$-equivariant quantizations with holomorphic-derivative vertices: for $\fsl_N$ and $\fgl_N$ ($\kappa=\tr$), $N\ge3$, order-$\hbar$ terms cancel the two-loop obstruction and the order-$\hbar^3$ obstruction vanishes for every such term ($\fsl_N$: the order-$\hbar^2$ and $\hbar^3$ terms are unique); for every simple $\g\ne\fsl_2$ with $\dim(\g^{*\otimes3})^\g=1$ (computed: $\fso_N$, $5\le N\le40$, $N\ne6$; $\fsp_N$, $4\le N\le40$; $G_2$, $F_4$, $E_6$, $E_7$, $E_8$) no order-$\hbar$ term of any bidegree cancels it. Order $\hbar^4$ is not analysed.
\end{abstract}

\tableofcontents

%======================================================================
\section{Setting and conventions}\label{sec-setting}
%======================================================================

\paragraph{Status of statements.} Theorem, Proposition, Lemma, Corollary: proved in the text, or established by the computational certificates of Appendix~\ref{app-cert} cited after the statement; each of these was recomputed independently of its first computation. Derived statement: a derivation is given, no independent recomputation was made. A statement read in a source carries the source and its section, equation or theorem number. Sources: S1 $=$ Li--Liu--Chang, arXiv:2611.24160 \cite{S1}; the other sources are listed in the bibliography.

\subsection{Network variables and reduced kernel}\label{sec-net}

Graph $\Gamma=(V,I)$, edges $e=(e(0),e(1))$ oriented and totally ordered (edge order), $n=|I|-|V|+1$, root $v_*$ with $\lambda_{v_*}=-\sum_{v\ne v_*}\lambda_v$, $\lambda_v\in\mathbb C^H$; $B_{ev}=\delta_{v,e(0)}-\delta_{v,e(1)}$ ($v\ne v_*$); cycle matrix $C\in\mathbb Z^{n\times|I|}$; $D=H+T$.
\begin{gather}
B^{\mathsf T}y=\lambda,\qquad \sum_eC_{re}t_ey_e=0,\qquad M=C\,t\,C^{\mathsf T},\qquad U=\det M,\qquad \psi^i_r=\sum_eC_{re}y^i_e\,dt_e,\\
\rho_n(M)=(\det M)^{-1/2}\Pf\big(-\tfrac12\,dM\,M^{-1}dM\big),\qquad \Pf(A)=\sum_{j=2}^n(-1)^jA_{1j}\Pf(A_{\hat1\hat\jmath}),\qquad\rho_{2k+1}=0,\label{eq-rho}\\
K_\Gamma(\lambda;z)=\int_{\RPp{|I|-1}}e^{-\sum_ey_e\cdot z_e}\,U^{-H}\rho_n(M)^T\wedge\prod_{i=1}^H\prod_{r=1}^n\psi^i_r .\label{eq-K}
\end{gather}
$\rho_{2k+1}=0$: Pfaffian of an antisymmetric matrix of odd order. $n=2$: $\Pf(A)=A_{12}$, $M^{-1}=(\det M)^{-1}\big(\begin{smallmatrix}M_{22}&-M_{12}\\-M_{12}&M_{11}\end{smallmatrix}\big)$, $dM_{12}\wedge dM_{12}=0$:
\begin{gather}
\big(dM\,{\rm adj}(M)\,dM\big)_{12}=M_{22}\,dM_{11}dM_{12}-M_{12}\,dM_{11}dM_{22}+M_{11}\,dM_{12}dM_{22},\notag\\
\rho_2=-\frac{M_{11}\,dM_{12}\,dM_{22}-M_{12}\,dM_{11}\,dM_{22}+M_{22}\,dM_{11}\,dM_{12}}{2(\det M)^{3/2}} .\label{eq-rho2}
\end{gather}
Orientation of $\RPp{|I|-1}$: $\iota_E\prod_edt_e$ in the edge order, $E=\sum_et_e\partial_{t_e}$. Factor order: $\rho_n$, then $i=1..H$ (outer), $r=1..n$ (inner). Degree: $n(H+T)=|I|-1$, hence
\begin{equation}
|I|=Dn+1,\qquad |V|=(D-1)n+2 .
\end{equation}
Density: $U^{-H}\rho_n^T\wedge\prod\psi=D(t)\,\iota_E(dt_0\wedge\dots\wedge dt_{|I|-1})$, $K_\Gamma=\int_{t>0,\sum t=1}D(t)\,e^{-\sum y\cdot z}\,dt_0\cdots dt_{|I|-2}$; $-\partial/\partial z^i_e$ inserts $y^i_e$. Determinants: $[u_1\cdots u_H]=\det(u_1,\dots,u_H)$, at $H=2$: $[ab]=a^1b^2-a^2b^1$.

Cycle bases. $C_T$: fundamental cycles of a spanning tree $T$, $(C_T)_{rc_r}=+1$ on the chord $c_r$; $\sigma_T=\varepsilon_T\det B_T$ ($\varepsilon_T$: sign of the permutation moving the tree edges in front). Positive basis: $\det\binom{B^{\mathsf T}}{C}>0$ ($=\#$ spanning trees); $K_\htk$ denotes $K_\Gamma$ in the positive basis.

\begin{proposition}[transformation rules]\label{prop-rules}
$C\to AC$, $|\det A|=1$: $K_\Gamma\to(\det A)^{D}K_\Gamma$. For $K_\htk$: reversal of $k$ edges with $z_e\to-z_e$: $(-1)^{Dk}$; permutation $\pi$ of the edge order: $\sgn(\pi)^{D+1}$; relabelling $\pi$ of the vertices (root $=$ last label): $\sgn(\pi)^D$. $T=0$, fundamental cycles of $T$, chords $c_r$:
\begin{equation}
K_\htk=\sigma_T^{\,H}\,\epsilon(n,H)\int e^{-\sum y\cdot z}\bigwedge_rd^Hy_{c_r},\qquad \epsilon(n,H)=(-1)^{nH+\binom n2\binom H2}.
\end{equation}
\end{proposition}
Under $C\to AC$: $M\to AMA^{\mathsf T}$, $\prod_r\psi^i_r\to\det A\prod_r\psi^i_r$, $\rho_n\to\det A\,\rho_n$.
\cert{C1}

\subsection{\texorpdfstring{Classical theory on $\mathbb R\times\mathbb C^2$}{Classical theory on R x C2}}\label{sec-classical}

Wang--Williams 2407.08667 v3 \cite{WW}, (8), (9):
\begin{gather}
X=\mathbb R\times\mathbb C^2\ni(\tau,z_1,z_2),\quad \Omega=dz_1\wedge dz_2,\quad \mathcal E=\Omega^*(\mathbb R)\otimes\Omega^{0,*}(\mathbb C^2)\otimes\g[1],\quad A=c+A^{(1)}+A^{(2)}+A^{(3)},\\
S_{\rm cl}(A)=\int_X\Omega\wedge\Big(\tfrac12\kappa(A,DA)+\tfrac16\kappa(A,[A,A])\Big),\qquad \omega(a,b)=\int_X\Omega\,\kappa(a,b),\quad \deg\omega=2-D=-1,\\
I=\int_X\Omega\,\tfrac16\kappa(A,[A,A]),\qquad D=d\tau\,\partial_\tau+d\zb_i\,\partial_{\zb_i}.
\end{gather}
$\kappa$: nondegenerate invariant form; $A^{(k)}$: form degree $k$, ghost number $1-k$; $c$ ghost; $A^{(2)},A^{(3)}$ antifields. $f_{abc}=\kappa([T_a,T_b],T_c)$.

Scaling. $\phi_l(\tau,z_1,z_2)=(\tau,l_1z_1,l_2z_2)$, $l\in(\mathbb C^\times)^2$; $\phi_g(\tau,z)=(\tau,gz)$, $g\in GL_2(\mathbb C)$:
\begin{equation}
\phi_g^*\Omega=\det(g)\,\Omega,\qquad S_{\rm cl}(\phi_g^*A)=\det(g)^{-1}S_{\rm cl}(A),\qquad \omega(\phi_g^*a,\phi_g^*b)=\det(g)^{-1}\omega(a,b);
\end{equation}
$(S_{\rm cl}/\hbar,\omega/\hbar)$ is invariant under $A\to\phi_g^*A$, $\hbar\to\det(g)^{-1}\hbar$. For $F_P=\int\Omega\,P(\partial)W(A,..,A)$, $P$ of bidegree $(a,b)$ in the holomorphic derivatives: $F_P(\phi_l^*A)=l_1^{a-1}l_2^{b-1}F_P(A)$.

\subsection{\texorpdfstring{Heat-kernel gauge on $\mathbb R\times\mathbb C^2$}{Heat-kernel gauge on R x C2}}\label{sec-gauge}

\begin{proposition}[edge data]\label{prop-gauge}
$D^*=-\iota_{\partial_\tau}\partial_\tau-4\sum_i\iota_{\partial_{\zb_i}}\partial_{z_i}$, $[D,D^*]=-\partial_\tau^2-4\sum_i\partial_{z_i}\partial_{\zb_i}=-\Delta_{\mathbb R^5}$;
\begin{gather}
[K_t]_\Omega=\frac{e^{-(\tau^2+|z|^2)/4t}}{4\,(4\pi t)^{5/2}}\,d\tau\,d^2\zb,\qquad \int_X\Omega\wedge[K_t]_\Omega=1,\qquad [D^*K_t]_\Omega=t^{-1}\iota_{E_x}[K_t]_\Omega,\quad E_x=\tfrac\tau2\partial_\tau+\zb^i\partial_{\zb^i},\\
[P]_\Omega=\int_0^\infty dt\,[D^*K_t]_\Omega=\frac{3}{16\pi^2r^5}\Big[\frac\tau2\,d\zb_1\wedge d\zb_2-d\tau\wedge(\zb_1d\zb_2-\zb_2d\zb_1)\Big],\qquad r^2=\tau^2+|z|^2,\\
D[P]_\Omega=D^*[P]_\Omega=0\ (r>0),\qquad \int_{S^4_R}\Omega\wedge[P]_\Omega=1,\qquad [P]_\Omega(-x)=-[P]_\Omega(x),\\
\Theta_e=[K_{t_e}]_\Omega-dt_e\wedge[D^*K_{t_e}]_\Omega=\frac1{4\pi^{5/2}}\,e^{-\eta_e^2-y_e\cdot\Delta_e}\,d\eta_e\wedge dy^1_e\wedge dy^2_e,\qquad(d_t+D)\Theta_e=0,\\
y_e=\frac{\zb_{e(0)}-\zb_{e(1)}}{4t_e},\qquad \eta_e=\frac{\tau_{e(0)}-\tau_{e(1)}}{2\sqrt{t_e}},\qquad \Delta_e=z_{e(0)}-z_{e(1)}+z_e .
\end{gather}
$[P]_\Omega$ equals Ishtiaque--Jeong--Zhou 2408.15712 \cite{IJZ} (2.12) with $(x,z,t)\to(z_1,z_2,\tau)$, $\epsilon_1\to1$. Wang--Williams gauge: $D^*_{\rm WW}=\frac12D^*$, $H_{\rm WW}=\frac12H$ under $x=\sqrt2\,\tau$; WW (22) as printed has $\int\Omega\wedge H_{\rm WW}=4$, so $P_{\rm WW}=4P$ componentwise for the printed kernel.
\end{proposition}
$[P]_\Omega$: $\int_0^\infty t^{-1}e^{-r^2/4t}(4(4\pi t)^{5/2})^{-1}dt=\Gamma(\frac52)(4/r^2)^{5/2}/(128\pi^{5/2})=3/(16\pi^2r^5)$. $\Theta_e$: $d\tau\,d^2\zb=2\sqrt t\,(4t)^2d\eta\,d^2y$, $32/(4\cdot32\pi^{5/2})=1/(4\pi^{5/2})$.
\cert{C7}

\begin{dstatement}[$U(2)$-equivariance of the scheme]\label{ds-U2}
For $g\in U(2)$: $\partial_{z'}=(g^{-1})^{\mathsf T}\partial_z$, $\iota_{\partial_{\zb'}}=(\bar g^{-1})^{\mathsf T}\iota_{\partial_{\zb}}$, $d\zb'=\bar g\,d\zb$, $g^{\mathsf T}\bar g=1$, hence
\begin{equation}
\textstyle\sum_id\zb'_i\partial_{\zb'_i}=\sum_id\zb_i\partial_{\zb_i},\qquad \sum_i\iota_{\partial_{\zb'_i}}\partial_{z'_i}=\iota_{\partial_{\zb}}^{\mathsf T}(g^{\mathsf T}\bar g)^{-1}\partial_z=\sum_i\iota_{\partial_{\zb_i}}\partial_{z_i},\qquad |gz|^2=|z|^2 ;
\end{equation}
$D$, $D^*$, $[D,D^*]$, $e^{-t[D,D^*]}$ and $P_\varepsilon^L=\int_\varepsilon^Ldt\,D^*K_t$ commute with $\phi_g^*$. With the scaling of $S_{\rm cl}$ and $\omega$ (Section~\ref{sec-classical}), the heat-kernel scheme is equivariant under $U(2)$ with $\hbar\to\det(g)^{-1}\hbar$. The scheme is not equivariant under $(\mathbb C^\times)^2\not\subset U(2)$.
\end{dstatement}

\subsection{Dictionary to S1}\label{sec-dict}

{\footnotesize
\renewcommand{\arraystretch}{1.3}
\begin{longtable}{@{}>{\raggedright\arraybackslash}p{2.6cm}>{\raggedright\arraybackslash}p{6.4cm}>{\raggedright\arraybackslash}p{6.4cm}@{}}
\toprule
 & S1 ($\mathbb C^2$) & $\mathbb R\times\mathbb C^2$\\
\midrule\endhead
$(H,T)$, $D$ & $(2,0)$, $2$ & $(2,1)$, $3$\\
$\Omega$ & $2\,dx^{+\dot+}\wedge dx^{+\dot-}$ (S1 (3.12)) & $dz_1\wedge dz_2$\\
$\Omega\wedge d\tau^T\wedge d^H\bar x$ & $-8\,d^4x_E$ & $4\,d\tau\,d^4x$\\
$[K_t]_\Omega$, $\int\Omega\wedge K_t=1$ & $-\dfrac{e^{-|x|^2/2t}}{32\pi^2t^2}dx^{-\dot+}dx^{-\dot-}$ (S1 (3.39)) & $\dfrac{e^{-(\tau^2+|z|^2)/4t}}{4\,(4\pi t)^{5/2}}d\tau\,d^2\zb$\\
currents & $y_{e,+\dot a}=\frac{(x_{e(0)}-x_{e(1)})^-_{\dot a}}{2t_e}$ (S1 (3.72)) & $y_e=\frac{\zb_{e(0)}-\zb_{e(1)}}{4t_e}$, $\eta_e=\frac{\tau_{e(0)}-\tau_{e(1)}}{2\sqrt{t_e}}$\\
edge form & $-\frac1{8\pi^2}e^{-y_e\cdot(x_{e(0)}-x_{e(1)}+z_e)}d^2y_e$ (S1 (3.73)) & $\frac1{4\pi^{5/2}}e^{-\eta_e^2-y_e\cdot\Delta_e}\,d\eta_e\,d^2y_e$\\
heat-edge sum & $(-1)^{j-1}$, $P$ odd (S1 (3.60)) & no signs, $P$ even\\
raw integral & $(-8\pi^2)^{-n}I^{\rm S1}$ & $(-1)^{\frac{n(n-1)}2}(8\pi^{5/2})^{-n}K_\htk$ ($n=2$: Proposition~\ref{prop-raw})\\
source kernel & $I^{\rm S1}=(-1)^{\frac{n(n+1)}2}K_\htk$ ($I^{\rm S1}$: S1 (3.69)) & $K_\htk$\\
\bottomrule
\end{longtable}}

``Heat-edge sum'' refers to the sum over the edge carrying $K_\varepsilon$. The signs $(-1)^{j-1}$ in the derivation chain of Proposition~\ref{prop-raw} are the orientation signs of the UV facets $t_j=\varepsilon$ in the fibre integration, not heat-edge signs. Schwinger times: $t=\tfrac12t^{\rm S1}$; GKW 2403.13049 \cite{GKW}: $t^{\rm GKW}=2t^{\rm S1}=4t$, $s=\eta$. The kernels are projective in $t$.
\cert{C4, C6}

\subsection{Raw two-loop integral}

\begin{proposition}[raw two-loop integral]\label{prop-raw}
Let $G_\Gamma(\lambda)=\lim_{\varepsilon\to0}\int_{X^{|V|-1}}\big(\bigwedge_{v\neq v_*}\Omega_v\big)e^{\sum_v\lambda_v\cdot z_v}\,D\big(\bigwedge_e[P_e]_\Omega\big)\big|_{\rm UV}$, $P_e=P(x_{e(0)}-x_{e(1)})$, $X_v$ oriented by $d\tau\,d^4x$, vertex order. For $\theta(1,3,3)$, $\theta(2,2,3)$ (engine labelling, Section~\ref{sec-graphs}) and the two thetas rooted at a trivalent vertex, at $(H,T)=(2,1)$:
\begin{equation}
G_\Gamma(\lambda)=-\frac{1}{64\pi^5}\,K_{\Gamma,\htk}(\lambda;0).
\end{equation}
\end{proposition}
Derivation chain: $\Theta_e$ odd, $\beta_e=[D^*K_{t_e}]$ even, $\Pi_*(dt_1\cdots dt_N\wedge a)=\int a\,d^Nt$;
\begin{gather}
\bigwedge_eP_e=(-1)^N\Pi_*\big(\textstyle\bigwedge_e\Theta_e\big),\qquad D\Pi_*=(-1)^N\Pi_*D,\qquad D\textstyle\bigwedge_e\Theta_e=-d_t\bigwedge_e\Theta_e,\\
D\textstyle\bigwedge_eP_e\big|_{\rm UV}=\sum_j(-1)^{j-1}\int_{t_j=\varepsilon}\big[\bigwedge_e\Theta_e\big]_{dt_1..\widehat{dt_j}..dt_N}=\int_{\partial{\rm UV},\ \iota_Edt^N}\bigwedge_e\Theta_e\quad(\text{no }(-1)^{\binom N2});
\end{gather}
extraction, reorderings, Gaussians, constant:
\begin{gather}
\theta^a_e=dX^a_e-c_aJ^a_edt_e,\quad c=(1,4,4),\quad J=(w,y^1,y^2):\qquad \varsigma^{\rm ext}=(-1)^{\binom{3n+1}2}=(-1)^{\frac{n(n-1)}2},\\
s_1=(-1)^{9n(|V|-1)}=(-1)^n,\qquad s_2=(-1)^{3\binom{|V|-1}2}=(-1)^n,\\
\int\prod_vd^4x_v\,e^{-\frac14z^\dagger\Lambda z+\lambda\cdot z}G(\zb)=(4\pi)^{2(|V|-1)}\Big(\frac{\prod_et_e}U\Big)^2G(4\Lambda^{-1}\lambda),\qquad \Lambda=B^{\mathsf T}t^{-1}B,\quad \det\Lambda=\frac{U}{\prod_et_e},\label{eq-holGauss}\\
\int\prod_vd\tau_v\,e^{-\frac14\tau^{\mathsf T}\Lambda\tau}=(4\pi)^{\frac{|V|-1}2}\Big(\frac{\prod_et_e}U\Big)^{1/2},\\ \langle w_ew_f\rangle=\tfrac12\big[\delta_{ef}/t_e-(C^{\mathsf T}M^{-1}C)_{ef}\big],\qquad U^{-1/2}\Big\langle\prod_r\psi^\tau_r\Big\rangle=\rho_n,
\end{gather}
\begin{equation}
(4\pi^{5/2})^{-|I|}(4(4\pi)^{5/2})^{|V|-1}4^{2n}\prod_e\frac{t_e^{5/2}}{2\sqrt{t_e}(4t_e)^2}=(8\pi^{5/2})^{-n} .
\end{equation}
\cert{C8, C9}

\begin{dstatement}[general $n$]\label{ds-gen-n}
$G_\Gamma=(-1)^{n(n-1)/2}(8\pi^{5/2})^{-n}K_{\Gamma,\htk}$ for every tight graph at $(H,T)=(2,1)$ (exponent and constant from the chain above for all $n$; pointwise identity computed only at $n=2$). General $(H,T)$: $G=\varsigma^{\rm ext}s_1s_2\,c_{H,T}^{-n}K$, $c_{H,1}=2v_H\pi^{H+1/2}$, $c_{H,0}=v_H\pi^H$, $\Omega\wedge d\tau^T\wedge d^H\zb=v_H\,d^{2H+T}x$.
\end{dstatement}

%======================================================================
\section{\texorpdfstring{Two-loop theta kernels at $T=1$ and $T=0$}{Two-loop theta kernels at T=1 and T=0}}\label{sec-theta}
%======================================================================

\subsection{Labelling and orientation}\label{sec-theta-lab}
$\theta(p,q,r)$: trivalent vertices $N=0$ and $N'=v_*$; paths $P\in\{A,B,C\}$ with $n_P=p,q,r$ edges $P_1,\dots,P_{n_P}$, all oriented $N\to N'$; edge order $(A_1..A_p,B_1..B_q,C_1..C_r)$; interior vertex $P_k$ between $P_k$ and $P_{k+1}$, numbered $1,2,\dots$ along $A$, then $B$, then $C$; source $\lambda^P_k$; $\Lambda_P=(\lambda^P_1,\dots,\lambda^P_{n_P-1})\in\mathbb C^{H\times(n_P-1)}$. Tight: $p+q+r=2D+1$, $n_P\le D$. Cycle matrix $C_3$: rows $A-B$, $B-C$.
\begin{gather}
R_P=\sum_jt_{P_j},\quad x^P_j=\frac{t_{P_j}}{R_P},\quad \sigma^P_k=\sum_{j>k}x^P_j,\quad w_P=\frac{R_P^{-1}}{\sum_{P'}R_{P'}^{-1}},\quad U=R_AR_B+R_BR_C+R_CR_A,\\
Q^P_j=\sum_{k<j}\lambda^P_k,\qquad \hat m_P=\sum_k\sigma^P_k\lambda^P_k,\qquad L=\lambda_N+\sum_P\hat m_P,\qquad y_{P_j}=w_PL-\hat m_P+Q^P_j,\\
D_\theta=\det\begin{pmatrix}\Lambda_A&\Lambda_B&0\\0&\Lambda_B&\Lambda_C\end{pmatrix}\ (T=1),\qquad
Q(v)=\det\begin{pmatrix}v&0&\Lambda_A&\Lambda_B&0\\0&v&0&\Lambda_B&\Lambda_C\end{pmatrix}\ (T=0).
\end{gather}
Both are $2H\times2H$ determinants.

\subsection{Densities}

\begin{theorem}[Schwinger-space densities]\label{thm-dens}
In the cycle basis $C_3$, edge order $A,B,C$:
\begin{align}
T=1:&\qquad D(t)=(-1)^{\frac{H(H+1)}2}\,\frac{D_\theta}{2\,U^{H+3/2}},\\
T=0:&\qquad D(t)=(-1)^{\frac{H(H+1)}2+1}\,\frac{R_AR_BR_C\,Q(L(t))}{U^{H+2}},\qquad L(t)=\lambda_N+\sum_{P,k}\sigma^P_k(t)\lambda^P_k .
\end{align}
At $T=1$ the density depends on $t$ only through $U$ and does not depend on $\lambda_N$. For any other labelling, orientation or edge order, $D_G=\sgn(\pi_{\rm std})D_{\rm std}$, $\pi_{\rm std}$ listing the edges in the order $A,B,C$.
\end{theorem}
Derivation chain:
\begin{gather}
\phi^i_P=\sum_jy^i_{P_j}dt_{P_j}=V^i\,d\log R_P+R_P\,d\hat m^i_P,\quad V=\frac{L}{\sum_PR_P^{-1}},\quad \psi_1=\phi_A-\phi_B,\ \ \psi_2=\phi_B-\phi_C,\\
a=d\log\frac{R_A}{R_B},\quad b=d\log\frac{R_B}{R_C},\quad a\wedge b=\frac{\Omega_R}{R_AR_BR_C},\quad \Omega_R=R_A\,dR_BdR_C-R_B\,dR_AdR_C+R_C\,dR_AdR_B,\\
\Big[\prod_i\psi^i_1\psi^i_2\Big]_{\rm fibre}=(-1)^{\frac{H(H-1)}2+H+q-1}\prod_PR_P^{n_P-1}D_\theta\,d\sigma,\qquad
\iota_E{\rm vol}=(-1)^{q-1}\prod_PR_P^{n_P-1}\,\Omega_R\wedge dx,\\
\rho_2=\frac{\Omega_R}{2U^{3/2}}\quad(M_{11}=R_A+R_B,\ M_{12}=-R_B,\ M_{22}=R_B+R_C;\ \eqref{eq-rho2}),\qquad R_P=1/w_P:\ U=\frac1{w_Aw_Bw_C}.
\end{gather}
$\Omega_R\wedge a=\Omega_R\wedge b=0$ removes the non-fibre terms at $T=1$. The signs of the intermediate steps are fixed by the final densities, which were computed pointwise.
\cert{C2}

\subsection{Integral representations}

\begin{theorem}[general sources and displacements]\label{thm-intrep}
With $d^2w=dw_A\,dw_B$ on $\Delta^2$ and $dx=\prod_Pdx^P_2\cdots dx^P_{n_P}$ on $\prod_P\Delta^{n_P-1}$ (Lebesgue):
\begin{align}
K^{T=1}[C_3](\lambda;z)&=(-1)^{\frac{H(H+1)}2}\frac{D_\theta}2\int_{\Delta^2}d^2w\prod_Pw_P^{H+\frac12-n_P}\int dx\,\exp\Big(\sum_v\lambda_v\cdot X_v\Big),\\
K^{T=0}[C_3](\lambda;z)&=(-1)^{\frac{H(H+1)}2+1}\int_{\Delta^2}d^2w\prod_Pw_P^{H-n_P}\int dx\,Q(L(x))\exp\Big(\sum_v\lambda_v\cdot X_v\Big),
\end{align}
\begin{equation}
Z_P=\sum_jz_{P_j},\quad W=\sum_Pw_PZ_P,\quad X_N=-W,\quad X_{P_k}=\sigma^P_k(Z_P-W)-\sum_{j>k}z_{P_j},\quad X_{N'}=0,
\end{equation}
and for a root at an interior vertex $\varrho$: $X^{(\varrho)}_v=X_v-X_\varrho$.
\end{theorem}
\cert{C3}

\subsection{\texorpdfstring{Values at $z=0$}{Values at z=0}}

\begin{theorem}[closed forms]\label{thm-closed}
\begin{align}
T=1:&\quad K[C_3](\lambda;0)=(-1)^{\frac{H(H+1)}2}\,\frac{D_\theta}2\,\frac{\prod_P\Gamma(H+\frac32-n_P)}{\Gamma(H+\frac32)\prod_P(n_P-1)!},\label{eq-T1}\\
T=0:&\quad K[C_3](\lambda;0)=(-1)^{\frac{H(H+1)}2+1}\,c_0\,Q(\bar L),\qquad c_0=\frac{\prod_P(H-n_P)!}{(H+1)!\prod_P(n_P-1)!},\label{eq-T0}\\
&\quad \bar L=\sum_v\mu_v\lambda_v,\qquad \mu_N=1,\quad\mu_{N'}=0,\quad \mu_{P_k}=\frac{n_P-k}{n_P}.\notag
\end{align}
Lemma used at $T=0$: the polarization of $Q$ vanishes on ${\rm span}(\Lambda_P)\times{\rm span}(\Lambda_P)$ for each $P$, so $\langle Q(L)\rangle=Q(\langle L\rangle)$, $\langle\sigma^P_k\rangle=(n_P-k)/n_P$. Basis: $\sgn\det\binom{B^{\mathsf T}}{C_3}=(-1)^{q+1}$, hence
\begin{equation}
K_\htk=(-1)^{(q+1)D}K[C_3].
\end{equation}
Root at an interior vertex: $K[C_3]$, as a function of the full source vector ($\sum_v\lambda_v=0$), does not depend on the root; $K_\htk[\varrho\leftrightarrow N']=(-1)^DK_\htk$.
\end{theorem}
\cert{C4}

Coefficients $c_1=K[C_3]/(\pi\,(-1)^{H(H+1)/2}D_\theta)$ at $T=1$ and $c_0$ at $T=0$; $\sgn_\htk$ is the sign of $K_\htk$ relative to $c_1\pi D_\theta$, resp.\ $c_0Q(\bar L)$:
{\small
\begin{longtable}{@{}l>{\raggedright\arraybackslash}p{11.2cm}l@{}}
\toprule
$T=1$, $H$ & $(p,q,r)$: $c_1$ & $\sgn_\htk$\\
\midrule
0 & $(111)$: $1$ & $+$\\
1 & $(122)$: $\frac13$ & $-$\\
2 & $(133)$: $\frac1{20}$;\ \ $(223)$: $\frac1{30}$ & $-,+$\\
3 & $(144)$: $\frac1{252}$;\ \ $(234)$: $\frac1{420}$;\ \ $(333)$: $\frac1{840}$ & $+,+,+$\\
4 & $(155)$: $\frac1{5184}$;\ \ $(245)$: $\frac1{9072}$;\ \ $(335)$: $\frac1{10080}$;\ \ $(344)$: $\frac1{22680}$ & $+,-,+,-$\\
5 & $(166)$: $\frac1{158400}$;\ \ $(256)$: $\frac1{285120}$;\ \ $(346)$: $\frac1{332640}$;\ \ $(355)$: $\frac1{798336}$;\ \ $(445)$: $\frac1{997920}$ & $-,-,-,-,-$\\
\midrule
$T=0$, $H$ & $(p,q,r)$: $c_0$ & $\sgn_\htk$\\
\midrule
1 & $(111)$: $\frac12$ & $+$\\
2 & $(122)$: $\frac16$ & $+$\\
3 & $(133)$: $\frac1{48}$;\ \ $(223)$: $\frac1{48}$ & $-,+$\\
4 & $(144)$: $\frac1{720}$;\ \ $(234)$: $\frac1{720}$;\ \ $(333)$: $\frac1{960}$ & $-,-,-$\\
5 & $(155)$: $\frac1{17280}$;\ \ $(245)$: $\frac1{17280}$;\ \ $(335)$: $\frac1{17280}$;\ \ $(344)$: $\frac1{25920}$ & $+,-,+,-$\\
6 & $(166)$: $\frac1{604800}$;\ \ $(256)$: $\frac1{604800}$;\ \ $(346)$: $\frac1{604800}$;\ \ $(355)$: $\frac1{967680}$;\ \ $(445)$: $\frac1{1088640}$ & $+,+,+,+,+$\\
\bottomrule
\end{longtable}}
At $T=1$: $K_\htk=s\,c_1\pi D_\theta$, $s=(-1)^{\lceil H/2\rceil+(H+1)(q+1)}$. Examples (vertex names for sources):
\begin{gather}
(1,1):\ K_\htk(\theta_{122})=-\tfrac\pi3\lambda_{B_1}\lambda_{C_1};\\
(2,1):\ K_\htk(\theta_{133})=-\tfrac\pi{20}[B_1B_2][C_1C_2],\quad K_\htk(\theta_{223})=+\tfrac\pi{30}[A_1B_1][C_1C_2];\\
(3,1):\ K_\htk(\theta_{333})=\tfrac{\pi}{840}\big([A_1A_2B_1][B_2C_1C_2]-[A_1A_2B_2][B_1C_1C_2]\big);\\
(2,0):\ K_\htk(\theta_{122})=-\tfrac16[\bar LB_1][\bar LC_1],\ \bar L=N+\tfrac12B_1+\tfrac12C_1;\\
(3,0):\ K_\htk(\theta_{133})=-\tfrac1{48}[\bar LB_1B_2][\bar LC_1C_2],\ \bar L=N+\tfrac23B_1+\tfrac13B_2+\tfrac23C_1+\tfrac13C_2;\\
(3,0):\ K_\htk(\theta_{223})=+\tfrac1{48}[\bar LA_1B_1][\bar LC_1C_2],\ \bar L=N+\tfrac12A_1+\tfrac12B_1+\tfrac23C_1+\tfrac13C_2 .
\end{gather}
Laplace expansions ($B'\sqcup B''=\{1,..,q-1\}$):
\begin{gather}
D_\theta=\sum_{|B'|=H-p+1}\varepsilon_1(B')[\Lambda_A\Lambda_{B'}][\Lambda_{B''}\Lambda_C],\quad \varepsilon_1=(-1)^{\frac{H(H+1)}2+\frac{p(p-1)}2+(H-p+1)(p-1)+\sum_{b\in B'}b},\\
Q(v)=\sum_{|B'|=H-p}\varepsilon_0(B')[v\Lambda_A\Lambda_{B'}][v\Lambda_{B''}\Lambda_C],\quad \varepsilon_0=(-1)^{\frac{H(H+1)}2+\frac{(p+1)(p+2)}2+(H-p)(p+1)+\sum_{b\in B'}b}.
\end{gather}
\cert{C4, C5}

\begin{proposition}[S1 two-loop fan]\label{prop-sourcevals}
$\theta(1,2,2)$ at $(2,0)$, S1 vertex $i=$ engine vertex $i-1$: $I^{\rm S1}(\lambda;0)=-\frac1{24}[\lambda_1,\lambda_3+2\lambda_2][\lambda_3,\lambda_1+2\lambda_2]$, $=\det(\lambda_1,\lambda_3)^2/24$ at $\lambda_2=0$ (S1 (5.6)); it equals minus the right side of S1 (7.13) with $(v_{11},v_{21},A)=(1,3,2)$.
\end{proposition}
\cert{C6}

%======================================================================
\section{Local functionals and the cochain model}\label{sec-h1}
%======================================================================

\subsection{Jets, translation invariants, weights}\label{sec-jets}

$U=\{\tau,\bar1,\bar2\}$, $x_\tau=\tau$, $x_{\bar i}=\zb_i$; $\Lambda=\Lambda[dx_v:v\in U]$, $dx_T$ ($T\subseteq U$) ordered $\tau<\bar1<\bar2$; $L=\Omega^*(\mathbb R)\otimes\Omega^{0,*}(\mathbb C^2)\otimes\g$, $d_L=\sum_{v\in U}dx_v\partial_v$, $\mathcal E=L[1]$. Jets at $0$: $L_0=\mathbb C[[\tau,z,\zb]]\otimes\Lambda\otimes\g$, coordinates $c^a_{T,\alpha}$, $\alpha\in\mathbb N^5$ (orders in $\tau,z_1,z_2,\zb_1,\zb_2$), ${\rm gh}(c^a_{T,\alpha})=1-|T|$, superfield $A^a=\sum_Tc^a_{T,0}dx_T$; $M=C^*_{\rm red}(L_0)$.
\begin{gather}
Qc^a_{T,\alpha}=\partial^\alpha\big[(d_LA^a+\tfrac12f^a{}_{bc}A^bA^c)_T\big],\qquad D_uc^a_{T,\alpha}=c^a_{T,\alpha+e_u},\qquad \eta_vc^a_{T,\alpha}=(\iota_{\partial/\partial x_v}A^a)_{T,\alpha},\label{eq-jetops}\\
Q^2=0,\qquad[Q,D_u]=0,\qquad Q\eta_v+\eta_vQ=-D_v,\qquad \eta_v\eta_w+\eta_w\eta_v=0 .\label{eq-jetid}
\end{gather}
\begin{gather}
V:=M\Big/\sum_uD_uM\quad(\text{translation-invariant local functionals }\textstyle\int d{\rm vol}\,m),\notag\\
C_\g:=C^*_{\rm red}(\g[[z_1,z_2]])\ (\text{polynomials in }c^a_{\emptyset,(0,\alpha_1,\alpha_2,0,0)}),\qquad \Cg:=C_\g/(D_{z_1}C_\g+D_{z_2}C_\g),\notag\\
\Cg=\{\text{alternating }k\text{-cochains of }\g[[z_1,z_2]]\text{, polynomial in }p_1..p_k,\ \text{restricted to }\textstyle\sum_jp_j=0\},\label{eq-C0}\\
\psi_S(\phi):=\int d^2z\wedge dx_{U\smallsetminus S}\wedge\Phi(A),\qquad \Phi(A):=\phi|_{c^a_\alpha\to\partial^\alpha A^a},\qquad {\rm gh}\,\psi_S(\phi)=\deg\phi-|S|,\notag\\
\Vht:=\psi_U(\Cg)=\Big\{\int\Omega\,P(\partial_z^{(1)},..,\partial_z^{(k)})\,W(A,..,A)\Big\},\notag\\
\Veta:=\{F\in V:\ \eta_vF=0\text{ in }V,\ v\in U\},\qquad \eta_U:=\eta_\tau\eta_{\bar1}\eta_{\bar2}.\notag
\end{gather}
$V$ (constant-coefficient densities modulo constant-coefficient total derivatives) is the definition of translation-invariant local functionals used throughout.

Weights. $\Gamma=\mathbb R_{>0}\times(\mathbb C^\times)^2\ni\gamma=(r,l_1,l_2)$: $(\tau,z_1,z_2)\mapsto(r\tau,l_1z_1,l_2z_2)$. For a monomial $F=\int d{\rm vol}\prod_jc_{T_j,\alpha_j}$:
\begin{equation}
\begin{gathered}
F(\gamma^*A)=r^{w_\tau}\,l_1^{a-1}l_2^{b-1}\,\bar l_1^{\,w_{\bar1}}\bar l_2^{\,w_{\bar2}}\,F(A),\\
w_v=\sum_j\big(\alpha_{j,v}+[v\in T_j]\big)-1\ (v\in U),\qquad (a,b)=\Big(\sum_j\alpha_{j,z_1},\sum_j\alpha_{j,z_2}\Big).
\end{gathered}\label{eq-weights}
\end{equation}
$w_v\ge-1$; $\psi_S(\phi)$ with $\phi$ of bidegree $(a,b)$ has $w=-1_{U\smallsetminus S}$. $V_{w,(a,b)}$: span of the monomials of weight $w$ and bidegree $(a,b)$.

\subsection{Cohomology}

\begin{theorem}[translation-invariant local cohomology]\label{thm-A}
$\g$ a finite-dimensional Lie algebra.
(A1) $H^g(V)=\bigoplus_{S\subseteq U}\psi_SH^{g+|S|}(\Cg)$, each $\psi_S$ injective on cohomology.
(A2) $H(V_{w,(a,b)})=\psi_SH(\Cg)_{(a,b)}$ for $w=-1_{U\smallsetminus S}$; $H(V_{w,(a,b)})=0$ for $w\notin\{0,-1\}^3$.
(A3) Every class of $H(V)$ has a representative $\psi_S(\phi)$ (holomorphic derivatives only).
\end{theorem}

\begin{theorem}[$\eta$-invariant local cohomology]\label{thm-B}
$\mathfrak t=\mathbb C^{5|3}={\rm span}\{D_u\}\oplus{\rm span}\{\eta_v\}$, $[Q,\eta_v]=-D_v$.
(B1) $\Veta=\eta_UV$, $\Veta\subset\bigoplus_{w\ge0}V_w$, $\Veta\cap V_{w=0}=\Vht$.
(B2) $\psi_U:\Cg\to\Vht$ is bijective, lowers the degree by $3$ and is a chain map; $\Vht\subset\Veta$ is a quasi-isomorphism:
\begin{equation}
H^g(\Veta)=H^{g+3}(\Cg)=H^{g+3}\big(\mathbb C\,d^2z\otimes^{\mathbb L}_{\mathbb C[\partial_{z_1},\partial_{z_2}]}C^*_{\rm red}(\g[[z_1,z_2]])\big).
\end{equation}
(B3) $H(\Veta)\to H(V)$ is injective with image $\psi_UH(\Cg)=H(V_{w=0})$.
(B4) Imposing only $\eta_{\bar1},\eta_{\bar2}$: $H^g(V^{\eta_{\bar1},\eta_{\bar2}})=H^{g+3}(\Cg)\oplus H^{g+2}(\Cg)=H^{g+2}\big(\mathbb C\,d\tau\,d^2z\otimes^{\mathbb L}_{\mathbb C[\partial_\tau,\partial_z]}C^*_{\rm red}(\g[[z]])\big)$, $\partial_\tau$ acting by $0$.
\end{theorem}

\begin{theorem}[invariants]\label{thm-C}
$\g$ reductive, $G$ connected reductive, ${\rm Lie}\,G=\g$, acting by ${\rm Ad}$: $\Cg=\bigoplus_{(a,b)}\Cg{}_{,(a,b)}$, $\dim C^{0,n}_{\g,(a,b)}<\infty$, $H(\Cg)=H((\Cg)^G)$, $H(V)=H(V)^G$, $H(\Veta)=H(\Veta)^G$.
\end{theorem}

\begin{proof}[Proof of Theorems~\ref{thm-A}--\ref{thm-C}]
P1 (cited). Williams 1809.02661 \cite{Wil} Def.~2.5: $C^*_{\rm loc}(L)={\rm Dens}_X\otimes_{D_X}C^*_{\rm Lie,red}(JL)$; Prop.~2.8: a BV theory $(\mathcal E,\omega,S)$ is a local $L_\infty$ algebra $L=\mathcal E[1]$ with $(-3)$-shifted pairing, and $\{S_{\rm cl},\cdot\}$ is the CE differential of $JL$, at jets \eqref{eq-jetops} up to one global sign. Williams eq.~(12): a translation-invariant local functional is $\sum\int F(D_1\varphi,..,D_k\varphi)\,d^dx$ with constant-coefficient $D_i$; modulo constant-coefficient total derivatives this is $V$. Gwilliam--Williams 1810.06534 \cite{GW18} Def.~1.14, Lemma~A.7: the translation-invariant de Rham model is $\mathbb C[d\tau,dz,d\zb]\otimes C^*_{\rm red}(L_0)$ with differential $Q+\sum_udx_uD_u$, i.e.\ the Koszul complex $K(D;M)$ of $(D_u)$ on $M$, whose top-form quotient is $V$ (the proof of Lemma~A.7, translation-invariant sections of a trivial bundle are fixed by their value at $0$, applies with the factor $\mathbb R$).

P2. $M^{(k)}=(\mathbb C[p_1..p_k]\otimes W^{\otimes k})^{S_k,\pm}$, $p_j\in\mathbb C^5$, $W={\rm span}\{c^a_T\}$, $D_u=$ multiplication by $P_u=\sum_jp_{j,u}$; $U_0=\{\sum_jc_jp_j:\sum_jc_j=0\}$ is $S_k$-stable, $\mathbb C[p_1..p_k]=\mathbb C[P]\otimes{\rm Sym}(U_0)$, $S_k$ trivial on $\mathbb C[P]$ $\Rightarrow$ $M^{(k)}=\mathbb C[P_1..P_5]\otimes({\rm Sym}(U_0)\otimes W^{\otimes k})^{S_k,\pm}$ is free over $\mathbb C[P]$ $\Rightarrow$ $K(D;M)\to V$ and $K(D_z;C_\g)\to\Cg$ are quasi-isomorphisms, and $M^{(k)}/(P)$ is \eqref{eq-C0}.

P3. ${\rm cost}(c^a_{T,\alpha})=(\alpha_\tau+[\tau\in T],\alpha_{\bar1}+[\bar1\in T],\alpha_{\bar2}+[\bar2\in T],\alpha_{z_1},\alpha_{z_2})$; monomials of $V_{w,(a,b)}$ have total cost $(w+1,a,b)$; the cost-$0$ generators $c^a_{\emptyset,0}$ are odd $\Rightarrow$ arity $\le\dim\g+|w+1|+a+b$, $\dim M_{w,(a,b)}<\infty$. $Q:(w,(a,b),g)\to(w,(a,b),g+1)$; $D_u$ raises $(w,a,b)$ by $e_u$; $\eta_v:(w,g)\to(w+e_v,g-1)$.

P4. $i:\g[[z]]\to L_0$ ($\tau$-, $\zb$-independent $0$-forms), $\pi:L_0\to\g[[z]]$ ($x_U=0$, $dx_U=0$): dg Lie maps, $\pi i={\rm id}$, $D_zi=iD_z$, $D_vi=0$, $\iota_vi=0$. $L_0=\mathbb C[[z]]\otimes(\bigotimes_{v\in U}B_v)\otimes\g$, $B_v=\mathbb C[[x_v]]\oplus\mathbb C[[x_v]]dx_v$, weight-$n$ part $\mathbb Cx_v^n\oplus\mathbb Cx_v^{n-1}dx_v$, $d\,x_v^n=n\,x_v^{n-1}dx_v$ $\Rightarrow$ $H(B_v)=\mathbb C$, $H(L_0)=i(\g[[z]])$ (K\"unneth on finite-dimensional weight parts). $i^*:M\to C_\g$ is a quasi-isomorphism: in each slice (P3) the arity filtration is finite, $Q=Q_1+Q_2$, ${\rm gr}\,i^*={\rm Sym}^k((i^\vee)[-1])$ (characteristic $0$). $i^*D_z=D_zi^*$, $i^*D_v=0$, $i^*\eta_v=0$.

P5. $K(D;M)\to K(D;C_\g)$ (by $i^*$) is a quasi-isomorphism (iterated mapping cones). On $C_\g$: $D_v=0$ ($v\in U$), $K(D;C_\g)=\Lambda[\varepsilon_\tau,\varepsilon_{\bar1},\varepsilon_{\bar2}]\otimes K(D_z;C_\g)\simeq\Lambda[\varepsilon_U]\otimes\Cg$, $\deg\varepsilon=-1$. With P1, P2:
\begin{equation}
H^g(V)=H^g(K(D;M))=\bigoplus_{S\subseteq U}\varepsilon_S\otimes H^{g+|S|}(\Cg).\label{eq-A-sum}
\end{equation}

P6. $D_0:=Q+\sum_iD_{z_i}\partial/\partial\varepsilon_{z_i}$, $\mathbf D:=D_0+\sum_{v\in U}D_v\partial/\partial\varepsilon_v$, $\Theta:=-\sum_{v\in U}\eta_v\partial/\partial\varepsilon_v$:
\begin{equation}
[Q,\eta_v\partial_{\varepsilon_v}]=-D_v\partial_{\varepsilon_v},\qquad [D_{z_i}\partial_{\varepsilon_{z_i}},\eta_v\partial_{\varepsilon_v}]=0,\qquad [[D_0,\Theta],\Theta]=0\ \Longrightarrow\ e^{-\Theta}D_0e^{\Theta}=\mathbf D .
\end{equation}
$\phi\in Z^n(\Cg)$, $\Phi$ a Koszul cycle of $K(D_z;C_\g)$ with $\varepsilon_z^0$-component $\phi$, $y=\varepsilon_S\otimes\pi^*\Phi$ ($D_0y=0$): $\kappa_S:=e^{-\Theta}y$ satisfies $\mathbf D\kappa_S=0$, $i^*\kappa_S=\varepsilon_S\otimes\Phi$, augmentation $\pm\eta_S\pi^*\phi=\pm\psi_S(\phi)$ ($\eta_vA=\iota_vA$ on superfields, $[\Phi(A)]_{dx_\emptyset}=\pi^*\phi$) $\Rightarrow$ (A1), (A3). Chain map: $Q\Phi(A)=d_L\Phi(A)+(Q_{\rm hol}\phi)(A)$, $[d_L\Phi(A)]_{dx_S}=\sum_{v\in S}\pm D_v[\Phi(A)]_{dx_{S-v}}$, $Q_{\rm hol}\phi\in D_zC_\g$ for $\phi$ closed in $\Cg$, $\psi_SD_z=D_z\psi_S$ $\Rightarrow$ $Q\psi_S(\phi)\in DM$.

P7. $\gamma\in\Gamma$ preserves orientation; $S_{\rm cl}(\gamma^*A)=(l_1l_2)^{-1}S_{\rm cl}(A)$, $\omega(\gamma^*a,\gamma^*b)=(l_1l_2)^{-1}\omega(a,b)$ $\Rightarrow$ $Q$ preserves the characters \eqref{eq-weights}; the characters $r^{w_\tau}l^{(a-1,b-1)}\bar l^{(w_{\bar1},w_{\bar2})}$ are pairwise distinct $\Rightarrow$ $V=\bigoplus_{w,(a,b)}V_{w,(a,b)}$ as complexes; with (A1) and P3: (A2).

P8. (a) For fixed $(a,\alpha)$: ${\rm span}\{c^a_{T,\alpha}\}_T=\Lambda[\eta]c^a_{\emptyset,\alpha}$, free of rank $1$; tensor powers with the derivation action are free (shearing isomorphism $h\otimes n\mapsto\sum h_{(1)}\otimes h_{(2)}n$ of the Hopf algebra $\Lambda[\eta]$, $\eta$ primitive); $S_k$-(sign-)invariants are direct summands, free over the local algebra $\Lambda[\eta]$; $\otimes\,\mathbb C[p]/(P)$: free. Hence $\Veta=\bigcap_v\ker\eta_v=\eta_UV$, $\ker(\eta_U|_V)=\sum_v\eta_vV$, $\eta_U:V/\sum_v\eta_vV\to\Veta$ bijective of degree $-3$; $w\ge-1$ $\Rightarrow$ $\Veta\subset\bigoplus_{w\ge0}V_w$; $\Veta\cap V_0=\eta_UV_{-1_U}$, $V_{-1_U}=\Cg$, $\eta_U=\pm\psi_U$ on it (P6) $\Rightarrow$ (B1).
(b) $V/\sum_v\eta_vV=\mathbb C\otimes_{U(\mathfrak t)}M$, $U(\mathfrak t)=\mathbb C[D_\tau,D_z,D_{\zb}]\otimes\Lambda[\eta_U]$, $d\eta_v=-D_v$. $M^{(k)}$ with $Q_1$ is K-projective over $U(\mathfrak t)$ ($L_0^\vee[-1]=U(\mathfrak t)\otimes{\rm span}\{c^a_{\emptyset,0}\}$, $Q_1c^a_{\emptyset,0}=0$; tensor powers by shearing; $S_k$ direct summand). $p:U(\mathfrak t)\to\mathbb C[D_z]$ ($D_v,\eta_v\to0$) is a quasi-isomorphism of dg algebras ($\mathbb C[D_v]\oplus\mathbb C[D_v]\eta_v$, $d(f\eta_v)=-fD_v$: $H=\mathbb C$). ${\rm gr}_k$: $\mathbb C\otimes_{U(\mathfrak t)}M^{(k)}\xrightarrow{i^*}\mathbb C\otimes^{\mathbb L}_{U(\mathfrak t)}C_\g^{(k)}=\mathbb C\otimes^{\mathbb L}_{\mathbb C[D_z]}C^{(k)}_\g=\Cg{}^{(k)}$, quasi-isomorphisms $\Rightarrow$ $\rho:V/\sum_v\eta_vV\to\Cg$ is a quasi-isomorphism.
(c) $H^g(\Veta)=H^{g+3}(V/\sum_v\eta_vV)=H^{g+3}(\Cg)$; $\rho\circ\eta_U^{-1}\circ\psi_U=\pm{\rm id}$ $\Rightarrow$ (B2). (d) (A1) $\Rightarrow$ (B3). (e) $\mathfrak t'={\rm span}\{D_u\}\oplus{\rm span}\{\eta_{\bar1},\eta_{\bar2}\}$: $V^{\eta_{\bar1},\eta_{\bar2}}=\eta_{\bar1}\eta_{\bar2}V\cong\mathbb C\otimes_{U(\mathfrak t')}M$ (degree $-2$), $U(\mathfrak t')\simeq\mathbb C[D_\tau,D_z]$, $i^*$ is $\mathbb C[D_\tau,D_z]$-linear with $D_\tau\to0$, $\mathbb C\otimes^{\mathbb L}_{\mathbb C[D_\tau]}\mathbb C=\mathbb C\oplus\mathbb C[1]$ $\Rightarrow$ (B4).

P9. $X\in\g\subset\g[[z]]$ (momentum $0$), $\iota_X:=X^a\partial/\partial c^a_{\emptyset,0}$: $\iota_XD_u=D_u\iota_X$, $\iota_X\eta_v=\eta_v\iota_X$; $L_X:=Q\iota_X+\iota_XQ$ is the coadjoint action of $X$. $C^{0,n}_{\g,(a,b)}$ is a finite-dimensional $G$-module, $Q_{\rm hol}$ is $G$-equivariant, $G$ reductive: $\Cg=(\Cg)^G\oplus C'$, $H(C')$ is a sum of nontrivial isotypic parts on which $\g$ acts by $L_X=[Q,\iota_X]$, i.e.\ by $0$ $\Rightarrow$ $H(C')=0$. The same on $V$ and $\Veta$.
\end{proof}
\cert{C26, C27, C28, C33}

Sources. Williams \cite{Wil} Def.~2.28 (translation invariance together with $\eta^iI=0$), Lemma~2.29 (proof: one spectral-sequence argument), Cor.~2.30 ($(\mathrm{Def}\,\mathcal E_V)^{\mathbb C^{2d|d}}\simeq\mathbb C\,d^dz\otimes^{\mathbb L}_{\mathbb C[\partial_z]}\mathcal O_{\rm red}(V_0[[z]])[d]$, flatness over $U(\mathbb C^{2d|d})$ asserted) are the $T=0$ statements on $\mathbb C^d$ corresponding to (B1), (B2); P8 proves the flatness and adds the factor $B_\tau$ and $\eta_\tau$; shift $[d]\leftrightarrow[H+T]=[3]$. Gwilliam--Williams \cite{GW18} Lemma~A.8 ($\mathbb C[dz][2d]\otimes C^\#_{\rm red}(\g[[z]])$, $T=0$) corresponds to $\Veta\cap V_{w=0}=\Vht$; in the model $V$, $\Veta$ has in addition a positive-weight part, nonzero at the cochain level and acyclic. Costello \cite{C16} App.~B, first sentence of the proof of Lemma~B.0.2 (every holomorphically translation-invariant Lagrangian on $\mathbb R\times\mathbb C^2$ is equivalent to $\int dz_1dz_2\,\Phi(\alpha\wedge D_1\alpha\wedge\dots)$, $D_i\in\mathbb C[\partial_{z_1},\partial_{z_2}]$, no $\partial_t,\partial_{\zb},dt,d\zb$): stated there without proof, for $\fgl_N$ uniformly in $N$; it is (B2), here for every finite-dimensional $\g$ at $e=0$. Costello's complex $\mathbb C[dz_1,dz_2][5]$ has no $dt$ factor, i.e.\ it corresponds to imposing $\eta_\tau$ as well; with $\eta_{\bar1},\eta_{\bar2}$ only, (B4) has the additional summand $H^{g+2}(\Cg)$.

\subsection{Classes outside the cochain model}

\begin{proposition}[translation invariance alone]\label{prop-HCEfalse}
For every $\g$ with $\kappa\neq0$ there are closed, non-exact ghost-number-$0$ elements of $V$ outside $\Veta$:
\begin{gather}
\phi_i:=\kappa(\partial_{z_i}c,c),\quad d\phi_i(x_0,x_1,x_2)=f(X_0,X_1,X_2)(p_0+p_1+p_2)^i=0,\quad [\phi_i]\neq0\in H^2(\Cg)_{e_i},\\
G_\epsilon:=\int d^2z\wedge\big(d\zb_1\wedge\kappa(\partial_{z_2}A,A)-d\zb_2\wedge\kappa(\partial_{z_1}A,A)\big)=\pm\psi_{\{\tau,\bar2\}}(\phi_2)\pm\psi_{\{\tau,\bar1\}}(\phi_1),\quad w=(0,-1,0),(0,0,-1),\\
G_c(\phi_g^*A)=G_{c'}(A),\quad c'=\det(g)^{-1}g\,c\,g^{\mathsf T}\ (g\in U(2)),\qquad g\,\epsilon\,g^{\mathsf T}=\det(g)\,\epsilon\ \Rightarrow\ G_\epsilon\ U(2)\text{-invariant},\\
T_\lambda:=\int d\tau\,d^2z\,d^2\zb\,\lambda(A_\tau)=\psi_{\{\tau\}}(\lambda(c)),\quad \lambda\in(\g/[\g,\g])^*,\quad w=(0,-1,-1),\ U(2)\text{-invariant}.
\end{gather}
Hence the identification ``translation-invariant local functionals of ghost number $g$ $\leftrightarrow$ $H^{g+3}(\Cg)$'' fails without the $\eta$-conditions (Theorem~\ref{thm-A}: $H^g(V)=H^{g+3}\oplus3H^{g+2}\oplus3H^{g+1}\oplus H^g$ of $\Cg$), and ghost-number-$0$ classes of arity $1$ ($T_\lambda$) and $2$ ($G_\epsilon$) exist in $V$.
\end{proposition}
\cert{C29}

\begin{proposition}[characters]\label{prop-chars}
Let $I_k$ satisfy $I_k(\phi_l^*A)=(l_1l_2)^{k-1}I_k(A)$ and, under the joint dilation $(\tau,z)\to(r\tau,rz)$, $I_k\to r^{2k-2}I_k$. For $\psi_S(\phi)$, $\phi\in\Cg{}_{,(a,b)}$, \eqref{eq-weights} gives $a-1-w_{\bar1}=k-1$, $b-1-w_{\bar2}=k-1$, $w_\tau+a+b-2+w_{\bar1}+w_{\bar2}=2k-2$, hence $w_\tau+2w_{\bar1}+2w_{\bar2}=0$, $w=0$, $S=U$, $(a,b)=(k,k)$. Under $U(1)^2$ alone, $T_\lambda$ ($S=\{\tau\}$, $(0,0)$) and the components of $G_\epsilon$ ($S=\{\tau,\bar i\}$) have the character of $I_1$; their dilation characters are $r^{-4}$, $r^{-2}$.
\end{proposition}
Dilation equivariance of the heat-kernel quantization is not proved here.
\cert{C30}

\subsection{\texorpdfstring{$\eta$-equivariance of the heat-kernel scheme}{eta-equivariance of the heat-kernel scheme}}\label{sec-eta}

\begin{proposition}[$\eta$-equivariance]\label{prop-eta}
$\eta_vI=\frac16\int\Omega\wedge\iota_v\kappa(A,[A,A])=0$ ($[\iota_v\beta]_{3\text{-form}}=\iota_v[\beta]_{4\text{-form}}$, $\Lambda$ has no $4$-forms); $\omega(\iota_va,b)+(-1)^{|a|}\omega(a,\iota_vb)=\int\Omega\wedge\iota_v\kappa(a,b)=0$;
\begin{equation}
[D^*,\iota_v]=0,\qquad [D,\iota_v]=\partial_v,\qquad [[D,D^*],\iota_v]=0\quad(v\in U);
\end{equation}
$e^{-t[D,D^*]}$ and $P^L_\varepsilon$ commute with $\iota_v$; $[K_t]_\Omega$ and $[P]_\Omega$ evaluated at $x-y$ are forms in $d(x_u-y_u)$ with coefficients functions of $x-y$, annihilated by the derivation $\iota^x_v+\iota^y_v$ of $\Omega^*(X\times X)$. Hence the graph weights $w_\Gamma(P,I)$ and $w_\Gamma(K_\varepsilon\text{ on }e_j;P\text{ elsewhere})$ are $\eta$-equivariant and $I[L]$, $O[L]$, $O=\lim_{L\to0}O[L]$ (Section~\ref{sec-reg}) are $\eta$-invariant.
\end{proposition}
The passage from the kernels to $I[L]$ uses $[\mathcal X,W(P,\cdot)]=0$ for $\mathcal X$ symplectic with $\mathcal X P=0$, $\mathcal XI=0$; the cut-offs $\varepsilon,L$ and the limit $L\to0$ are not written out. No Feynman-graph computation of $\eta_vO[L]$ was made.
\cert{C31}

\begin{dstatement}[vertices in $\Vht$]\label{ds-Vht-eta}
For every $F=\int\Omega\,\Phi(A)\in\Vht$: $[\iota_v\Phi(A)]_{3\text{-form}}=\iota_v[\Phi(A)]_{4\text{-form}}=0$, so $\eta_vF=0$ identically. With Proposition~\ref{prop-eta}, the scheme with interaction $I+\sum_{k\ge1}\hbar^kI_k$, $I_k\in\Vht$, is $\eta$-equivariant; its obstructions $O_n$ are $\eta$-invariant, Corollary~\ref{cor-O2} applies to them, and their terms with $\partial_\tau$ or $\partial_{\zb}$ on external fields are exact.
\end{dstatement}

\begin{corollary}[obstruction classes]\label{cor-O2}
Let $O_n\in V^1$ be closed and $\eta$-invariant. Then $O_n=\psi_U(o)+\sum_{w\ne0}O_n^{(w)}$ with $O_n^{(w)}$ exact in $\Veta$ (Theorem~\ref{thm-B}); the evaluation $j^*$ on fields with $\tau$- and $\zb$-independent components $Xe^{p\cdot z}\otimes dx_T$ annihilates every term with $\partial_\tau$ or $\partial_{\zb}$, is defined on $V$ (every term of $D_vm$, $v\in U$, contains a $\partial_v$), and $o\mapsto\psi_U(o)\bmod D_z$ is injective; hence $o$ is fixed by the evaluation of $O_n$ on holomorphic plane waves, and
\begin{equation}
[O_n]=0\text{ in }H^1(V)\iff[O_n]=0\text{ in }H^1(\Veta)\iff[o]=0\text{ in }H^4(\Cg)\iff[o]=0\text{ in }H^4((\Cg)^G).
\end{equation}
Terms of $O_n$ with $\partial_\tau$ or $\partial_{\zb}$ on external fields have $w\neq0$ and are exact.
\end{corollary}
\cert{C27, C32}

\begin{proposition}[order-$\hbar$ terms in $\Veta$]\label{prop-J1eta}
Every closed $J\in V^{\eta,0}$ is $J=\psi_U(j)+QK$, $j\in Z^3(\Cg)$, $K\in V^{\eta,-1}$; $\psi_U(j)$ has $U(1)^2$-character $(0,0)$ iff $j$ has bidegree $(1,1)$.
\end{proposition}
\cert{C32}

\begin{dstatement}[representative independence]\label{ds-homotopy}
At order $\hbar$, $Q_0Q_{(1)}+Q_{(1)}Q_0=0$; for $Q_0J_1=0$, $J_1\to J_1+Q_0K$ changes $Q_{(1)}J_1+\frac12\{J_1,J_1\}$ by $Q_0\big(-Q_{(1)}K+\{K,J_1\}+\frac12\{K,Q_0K\}\big)$. ($Q_0$, $Q_{(1)}$: Section~\ref{sec-qme}.) Not checked at finite scale $L$. The one-loop term $Q_{(1)}J$ for $J\notin\Vht$ is not computed.
\end{dstatement}

%======================================================================
\section{\texorpdfstring{The two-loop obstruction of commutative CS on $\mathbb R\times\mathbb C^2$}{The two-loop obstruction of commutative CS on R x C2}}\label{sec-obs}
%======================================================================

\subsection{Regularization and obstruction}\label{sec-reg}

WW \cite{WW} (26)--(30), Def.~2.5, Cor.~2.8, Sec.~3.3.2:
\begin{gather}
P_\varepsilon^L=\int_\varepsilon^Ldt\,D^*K_t,\quad DP_\varepsilon^L=K_\varepsilon-K_L,\quad I[L]=\lim_{\varepsilon\to0}W(P_\varepsilon^L,I),\quad W(P,I)=\sum_\gamma|\Aut\gamma|^{-1}\hbar^{g(\gamma)}w_\gamma(P,I),\\
O[L]=\hbar\,e^{-I[L]/\hbar}(Q+\hbar\Delta_L)e^{I[L]/\hbar}=QI[L]+\hbar\Delta_LI[L]+\tfrac12(I[L],I[L])_L,\\
O[L]=\lim_{\varepsilon\to0}\sum_\Gamma|\Aut\Gamma|^{-1}\sum_{j\in I(\Gamma)}w_\Gamma\big(K_\varepsilon\text{ on }e_j;\ P^L_\varepsilon\text{ on }e\neq e_j\big),\qquad O=\lim_{L\to0}O[L].
\end{gather}
No counterterms (WW Thm~1.1, UV finiteness). $O_n$: the $b_1=n$ part of $O$; it carries $\hbar^n$.

\begin{dstatement}[closedness of the leading obstruction]\label{ds-closed}
$(Q+\hbar\Delta_L)^2=0$, $O[L]$ odd, $O[L]^2=0$:
\begin{equation}
0=(Q+\hbar\Delta_L)\big(O[L]e^{I[L]/\hbar}\big)=\big(QO[L]+\hbar\Delta_LO[L]+(I[L],O[L])_L\big)e^{I[L]/\hbar}-\hbar^{-1}O[L]^2e^{I[L]/\hbar} .
\end{equation}
If $O[L]=\hbar^nO_n[L]+O(\hbar^{n+1})$, then $QO_n[L]+(I^{(0)}[L],O_n[L])_L=0$ and, as $L\to0$, $\{S_{\rm cl},O_n\}=0$; in cochains $d\,o_n=0$. (Sign conventions of $(\cdot,\cdot)_L$ not tracked; closedness does not depend on them.)
\end{dstatement}

Vanishing. (i) Odd $n$ at $T=1$: reflection $s(\tau,z)=(-\tau,z)$ commutes with $\partial_z$, $s^*[P]_\Omega=-[P]_\Omega$, $s^*[K_\varepsilon]_\Omega=-[K_\varepsilon]_\Omega$, orientation of $X^{|V|-1}$: $(-1)^{|V|-1}$, $e^{\lambda\cdot z}$ invariant $\Rightarrow$ $G_\Gamma=(-1)^{|I|+|V|-1}G_\Gamma=(-1)^{b_1}G_\Gamma$; equivalently $\rho_{2k+1}=0$ in \eqref{eq-rho}; WW Prop.~3.23. (ii) Self-loops: $[K_t]_\Omega$, $[P]_\Omega$ are forms in $d(x-y)$ and vanish on the diagonal; also $f_{abc}\kappa^{bc}=0$. (iii) Non-tight graphs vanish by form degree.

\begin{proposition}[counting]\label{prop-count}
One heat edge, cubic vertices, $E$ external fields $A$ (total degree $1$ each):
\begin{equation}
|I|=Dn+1,\quad |V|=(D-1)n+2,\quad E=3|V|-2|I|=(D-3)n+4,\quad {\rm gh}=E-D=(D-3)(n-1)+1 .
\end{equation}
$D=3$: $E=4$, ${\rm gh}=1$ for all $n$; tight graphs with degrees $\le3$: $n=1$ box, $n=2$ $\theta(1,3,3)$, $\theta(2,2,3)$ (degrees $[2,2,2,2,3,3]$). $D=2$: $(n,E,{\rm gh})=(1,3,1),(2,2,0),(3,1,-1),(4,0,-2)$. BF vertex $\langle B,[A,A]\rangle$ with $\langle AB\rangle$ propagators: $b_1\le1$.
\end{proposition}
At $D=3$, $n=2$: $3|V|=2|I|+4$, WW (74).
\cert{C11}

\subsection{Graphs and kernels}\label{sec-graphs}

Engine labelling ($N=0$, $N'=v_*=5$, external vertices $1,2,3,4$ in path order, edges in the order listed, positive cycle basis):
\begin{align}
\theta(1,3,3):&\ (0,5),(0,1),(1,2),(2,5),(0,3),(3,4),(4,5);\quad C=\begin{pmatrix}-1&1&1&1&0&0&0\\-1&0&0&0&1&1&1\end{pmatrix},\\
\theta(2,2,3):&\ (0,1),(1,5),(0,2),(2,5),(0,3),(3,4),(4,5);\quad C=\begin{pmatrix}1&1&-1&-1&0&0&0\\-1&-1&0&0&1&1&1\end{pmatrix}.
\end{align}
$|\Aut\theta(1,3,3)|=|\Aut\theta(2,2,3)|=4$ (external legs unlabelled). Engine versus $C_3$: $\theta(1,3,3)$ rows $(B-A,C-A)=\big(\begin{smallmatrix}-1&0\\-1&-1\end{smallmatrix}\big)C_3$, $\det=+1$; $\theta(2,2,3)$ rows $(A-B,C-A)=\big(\begin{smallmatrix}1&0\\-1&-1\end{smallmatrix}\big)C_3$, $\det=-1$; $K_{\rm eng}=(\det)^3K[C_3]$.

\begin{theorem}[obstruction kernels]\label{thm-kernels}
At $(H,T)=(2,1)$, in the engine labelling, with $p_v=\lambda_v$ ($v=1..4$) and $\lambda_N$ arbitrary:
\begin{equation}
K_{133}(\lambda;0)=-\frac{\pi}{20}[p_1p_2][p_3p_4],\qquad K_{223}(\lambda;0)=+\frac{\pi}{30}[p_1p_2][p_3p_4].
\end{equation}
Both are \eqref{eq-T1} with $K_\htk=s\,c_1\pi D_\theta$, $s=-1,+1$.
\end{theorem}
\cert{C12}

\subsection{Wick sign and assembly}\label{sec-wick}

Rule: vertex $f_{h_1h_2h_3}A_{h_1}A_{h_2}A_{h_3}$ in a chosen local half-edge order; $\langle A^a(x)A^b(y)\rangle=s_P\hbar\kappa^{ab}[P](x-y)$, $[P]$ even; Koszul rule for the total degree;
\begin{gather}
\sigma_\Gamma=\sgn(\text{source}\to\text{target}),\quad\text{source}=\text{half-edges in local order, vertices in vertex order},\notag\\
\text{target}=(e,0),(e,1)\ \text{for }e\text{ in edge order, then }x_{v_1},\dots,x_{v_4},\notag\\
w_\Gamma(A)=U\hbar^2\sigma_\Gamma\int_X\Omega\wedge\big[G_\Gamma(\partial^{(1)},..,\partial^{(4)})\,C_\Gamma^{a_1a_2a_3a_4}A_{a_1}A_{a_2}A_{a_3}A_{a_4}\big]_{3\text{-form}},\\
C_\Gamma=\prod_vf_{(\text{local order})}\prod_e\kappa^{-1}_{(e,0)(e,1)},
\end{gather}
$O_2=\sum_\Gamma|\Aut\Gamma|^{-1}w_\Gamma$; $U=\pm1$ collects $s_P^7$, $(-1)^{5\binom62}$, the sign of the heat-edge identity, the placement of the external $3$-form and the global sign between $\{S_{\rm cl},\cdot\}$ and $Q$; it is one sign for all graphs with $(|V|,|I|,E)=(6,7,4)$. $U$ is not determined here.

\begin{proposition}[relative Wick sign]\label{prop-wick}
In the engine labelling with the local order $L'$ (externals first, then incident half-edges by (edge, end)): $\sigma_{133}=\sigma_{223}=-1=\sgn(L\to L')$ and $\sigma_\Gamma C_\Gamma=W^\Gamma$,
\begin{align}
W_{133}(A_1,..,A_4)&=-f_{xyz}\,f_{a_1y'u}\,f_{a_2u'w}\,f_{a_3z's}\,f_{a_4s'r}\,f_{x'w'r'},\\
W_{223}(A_1,..,A_4)&=-f_{xyz}\,f_{a_1x'u}\,f_{a_2y'w}\,f_{a_3z's}\,f_{a_4s'r}\,f_{u'w'r'},
\end{align}
primed pairs contracted with $\kappa^{-1}$. The relative sign of $\theta(1,3,3)$ and $\theta(2,2,3)$ in the $K_\Gamma W^\Gamma$ convention is $+1$; $\sigma_\Gamma K_\Gamma C_\Gamma$ is invariant under vertex relabelling (root included), edge reversal, edge order and local orders.
\end{proposition}
Scope. The rule $\sigma_\Gamma=\sgn(\text{source}\to\text{target})$ is derived from the Koszul convention (odd superfields, even contraction); it is not an end-to-end BV computation. The certificates compute $\sigma_\Gamma$, the identity $\sigma_\Gamma C_{L'}=W^\Gamma$ and the invariance. The relative sign enters the coefficient of $[O_2]$; the non-exactness of $O_2$ holds for all three choices $-3F_{133}\pm2F_{223}$, $3F_{133}+2F_{223}$ (Theorem~\ref{thm-allsimple}).
\cert{C10}

\begin{theorem}[assembly]\label{thm-assembly}
\begin{gather}
O_2=\cN\sum_\Gamma|\Aut\Gamma|^{-1}\int_X\Omega\,K_\Gamma(\partial;0)\,W^\Gamma(A,A,A,A),\qquad \cN=-\frac{U\hbar^2}{64\pi^5},\\
O_2=\frac{\pi\cN}{240}\int_X\Omega\,[\partial_1\partial_2][\partial_3\partial_4]\,W_{\rm tot}(A,A,A,A),\qquad W_{\rm tot}=-3W_{133}+2W_{223},\qquad \frac{\pi\cN}{240}=-\frac{U\hbar^2}{15360\,\pi^4},
\end{gather}
$\partial_i$: holomorphic derivative on the $i$-th $A$, $[\partial_i\partial_j]=\partial^{(i)}_{z_1}\partial^{(j)}_{z_2}-\partial^{(i)}_{z_2}\partial^{(j)}_{z_1}$; evaluated on holomorphic plane waves $e^{\lambda\cdot z}$. By Corollary~\ref{cor-O2} this determines the class of $O_2$ in $H^1(V)$.
\end{theorem}
Assembly chain: $|\Aut|=4$, $-\frac\pi{20}=\frac{\pi}{60}(-3)$, $\frac\pi{30}=\frac\pi{60}\cdot2$, $\frac{\cN}4\cdot\frac\pi{60}=\frac{\pi\cN}{240}$.
\cert{C10, C11}

\begin{theorem}[colour factors, $\kappa=\tr$]\label{thm-colour}
$\fgl_N$, $\tr(ijkl)=\tr(X_iX_jX_kX_l)$:
\begin{multline}
W_{\rm tot}=-10\tr(1234)-(10+3N^2)\tr(1243)+20\tr(1324)-(10+3N^2)\tr(1342)+20\tr(1423)-10\tr(1432)\\
-26N\tr(12)\tr(34)+4N\tr(13)\tr(24)+4N\tr(14)\tr(23)\\
+3N\big[\tr(1)\tr(234)+\tr(1)\tr(243)+\tr(123)\tr(4)+\tr(124)\tr(3)\\
+\tr(132)\tr(4)+\tr(134)\tr(2)+\tr(142)\tr(3)+\tr(143)\tr(2)\big]
+20\big[\tr(1)\tr(2)\tr(34)+\tr(12)\tr(3)\tr(4)\big]\\
-10\big[\tr(1)\tr(23)\tr(4)+\tr(1)\tr(24)\tr(3)+\tr(13)\tr(2)\tr(4)+\tr(14)\tr(2)\tr(3)\big].
\end{multline}
$\fsl_N$: the first nine terms; separately
\begin{align}
W_{133}&=2\tr(1234)+(2+N^2)\tr(1243)-4\tr(1324)+(2+N^2)\tr(1342)-4\tr(1423)+2\tr(1432)\notag\\&\quad+6N\tr(12)\tr(34),\\
W_{223}&=-2\tr(1234)-2\tr(1243)+4\tr(1324)-2\tr(1342)+4\tr(1423)-2\tr(1432)-4N\tr(12)\tr(34)\notag\\&\quad+2N\tr(13)\tr(24)+2N\tr(14)\tr(23).
\end{align}
$\fsl_2$: $W_{\rm tot}=-104\,k_{12}k_{34}+16\,k_{13}k_{24}+40\,k_{14}k_{23}$, $k_{ij}=\tr(X_iX_j)$.
\end{theorem}
\cert{C13}

\subsection{Cochains and closedness}\label{sec-cochains}

$x_i=X_ie^{p_i\cdot z}$, $[Xe^{p\cdot z},Ye^{q\cdot z}]=[X,Y]e^{(p+q)\cdot z}$, $(\sigma.f)(x_1..x_k)=f(x_{\sigma1}..x_{\sigma k})$, $\alt f=\sum_\sigma\sgn(\sigma)\,\sigma.f$;
\begin{gather}
(d\omega)(x_0,..,x_k)=\sum_{i<j}(-1)^{i+j}\,\omega([x_i,x_j],x_0,..,\widehat{x_i},..,\widehat{x_j},..,x_k),\\
F=\int\Omega\,P(\partial^{(1)},..,\partial^{(k)})\,W(A,..,A)\ \longleftrightarrow\ w=\alt(PW),\qquad F=\tfrac1{k!}w(A,..,A),\label{eq-dict}\\
F_G=\sum_{s\in S_4}\sgn(s)\,[p_{s1}p_{s2}][p_{s3}p_{s4}]\,W_G(X_{s1},..,X_{s4}),\\ \omega_{133}:=\alt\big(\tfrac{60}\pi K_{133}W_{133}\big)=-3F_{133},\quad \omega_{223}=2F_{223},\\
o_2:=\text{cochain of }O_2/\hbar^2=\frac{\pi\nu}{240}\big(-3F_{133}+2F_{223}\big)=\pi\nu\,o',\\ o'=\sum_G|\Aut G|^{-1}\alt(K_G/\pi\,W_G),\qquad \nu=\cN/\hbar^2=-\frac{U}{64\pi^5}.
\end{gather}
Ghost number $1$ $\leftrightarrow$ CE degree $4=1+H+T$ (Theorem~\ref{thm-B}); bidegree $(a,b)$ $=$ degree in the first, second momentum components; $O_2$ has bidegree $(2,2)$.

\begin{proposition}[closedness]\label{prop-closed}
$dF_{133}=dF_{223}=0$ for formal $\fgl_N$, $\fsl_N$ ($\kappa=\tr$); for $\fsl_2$, $\fgl_2$, $\fsl_3$, $\fgl_3$, $\fsl_4$, $\fgl_4$, $\fsl_5$, $\fgl_5$ ($\kappa=\tr$); for $\fso_5$, $\fsp_4$, $\fso_6$, $\fso_7$, $\fsp_6$, $\fso_8$, $G_2$, $F_4$, $\fso_9$, $\fsp_8$, $\fso_{12}$, $E_6$, $E_7$, $E_8$ (Killing form).
\end{proposition}
For $E_6,E_7,E_8$ also from $F_{223}=x\beta+aG_{kk}$, $dG_{kk}=0$ (Theorem~\ref{thm-simple}). For the other simple $\g$ closedness of $o_2$ follows from Derived statement~\ref{ds-closed}.
\cert{C14}

\subsection{\texorpdfstring{The class for $\fsl_N$ and $\fgl_N$}{The class for sl(N) and gl(N)}}\label{sec-class-slN}

\begin{theorem}[$H^4_{(2,2)}$ and {$[O_2]$}, $\kappa=\tr$]\label{thm-class-slN}
(i) Formal $\fgl_N$ (faithful for $N\ge5$): $\dim C^3=8$, $\dim C^4=36$, ${\rm rank}\,d_3=6$, ${\rm rank}\,d_4=26$, $\dim Z^4=10$, $\dim B^4=6$, $\dim H^4_{(2,2)}=4$, basis $\{[F_1],[F_0],[z_8],[z_9]\}$, $F_{223}=F_0+NF_1$ ($F_k$ $N$-independent),
\begin{equation}
[F_{223}]=N[F_1]+[F_0],\qquad [F_{133}]=-N[F_1]-[F_0],\qquad [F_{133}+F_{223}]=0 .
\end{equation}
Formal $\fsl_N$: $\dim C^3=3$, $\dim C^4=15$, ${\rm rank}\,d_3=3$, ${\rm rank}\,d_4=11$, $Z^4=4$, $B^4=3$, $H^4_{(2,2)}=\mathbb C[F_1]$, $[F_{223}]=N[F_1]=-[F_{133}]$.

(ii) Fixed $N$, $(\dim C^4,Z^4,B^4,H^4)$: $\fgl_1$ $(1,1,0,1)$; $\fgl_2$ $(22,8,5,3)$; $\fgl_3$ $(35,10,6,4)$; $\fgl_4,\fgl_5$ $(36,10,6,4)$; $\fsl_2$ $(5,2,2,0)$; $\fsl_3$ $(14,4,3,1)$; $\fsl_4,\fsl_5$ $(15,4,3,1)$. $[F_{223}]\neq0$ exactly for $N\ge3$; $[F_{133}+F_{223}]=0$ for all nine.

(iii)
\begin{equation}
[o_2]=\frac{\pi\nu}{240}(3+2)[F_{223}]=\frac{\pi\nu}{48}[F_{223}],\qquad [O_2]=\frac{\pi\cN}{48}[F_{223}]=-\frac{U\hbar^2}{3072\,\pi^4}[F_{223}],
\end{equation}
nonzero for $\fsl_N,\fgl_N$ with $N\ge3$, zero for $\fsl_2,\fgl_1,\fgl_2$.
\end{theorem}
Coefficient chain: $-3[F_{133}]+2[F_{223}]=5[F_{223}]$; $\frac{\pi}{240}\cdot5=\frac\pi{48}$; $\frac{\pi}{48}\cdot\frac{-U\hbar^2}{64\pi^5}=-\frac{U\hbar^2}{3072\pi^4}$. Cochain identities ($\kappa=\tr$, $\beta:=d([p_1p_2]^2\tr(1[23]))$): $F_{133}+F_{223}=2N^2\beta$, $o'=\frac1{48}F_{223}-\frac{N^2}{40}\beta$ ($\fsl_N$ and $\fgl_N$).
\cert{C15, C43}

\subsection{\texorpdfstring{Simple $\g$}{Simple g}}\label{sec-simple}

$G$: adjoint group; $SL_2(\mathbb C)$ acts on $(z_1,z_2)$; ${\rm inv}$: $G\times SL_2$-invariant part. $q_i$: momenta of a 3-cochain ($q_2=-q_0-q_1$), $f(X,Y,Z)=\kappa([X,Y],Z)$;
\begin{equation}
\begin{gathered}
\beta:=d\big([q_0q_1]^2f\big)=\sum_{i<j}(-1)^{i+j}[p_kp_l]^2\,\kappa([X_i,X_j],[X_k,X_l]),\qquad\{k<l\}=\{0,1,2,3\}\setminus\{i,j\},\\
G_{kk}:=\alt\big([p_0p_1][p_2p_3]\,\kappa_{01}\kappa_{23}\big),
\end{gathered}
\end{equation}
using $[p_i+p_j,p_k]=[p_kp_l]$ on $\sum p=0$; $\kappa_{ij}=\kappa(X_i,X_j)$.

\begin{lemma}[reduction]\label{lem-red}
(a) $w\in C^4_{\rm inv}$, $w=d\eta$ $\Rightarrow$ $w=d(R\eta)$, $R$ the Reynolds projector of $G\times SL_2$ on the 3-cochains of momentum degree 4; $H^4_{\rm inv}\hookrightarrow H^4_{(2,2)}$. (b) The $SL_2$-invariant quartics in $(q_0,q_1)$ are $\mathbb C[q_0q_1]^2$, $[q_{s0}q_{s1}]=\sgn(s)[q_0q_1]$; with $\dim(\Lambda^3\g^*)^\g=1$: $C^3_{\rm inv}=\mathbb C[q_0q_1]^2f$ and $B^4\cap C^4_{\rm inv}=\mathbb C\beta$. (c) For $X\in Z^4_{\rm inv}$: $[X]=0\iff X\in\mathbb C\beta$. (d) $\g$ simple, $T\in(\g^{*\otimes k})^\g$, $T(b,\cdot,..)=0$ for one $b\ne0$ $\Rightarrow$ $T=0$.
\end{lemma}
\cert{C16}

\begin{theorem}[classes for the computed simple algebras]\label{thm-simple}
$\kappa$ the Killing form; $\mathcal L=\{\fsl_3,..,\fsl_7,\ \fso_5=\fsp_4,\ \fso_7,\fsp_6,\fso_8,\fso_9,\fsp_8,\fso_{10},\fso_{11},\fsp_{10},\fso_{12},\ G_2,F_4,E_6,E_7,E_8\}$ (and $\fso_6=\fsl_4$). For $\g\in\mathcal L$:
\begin{equation}
F_{133}+F_{223}=\tfrac12\beta,\qquad F_{133},\ F_{223},\ -3F_{133}\pm2F_{223},\ G_{kk}\notin\mathbb C\beta,\qquad dG_{kk}=0,\qquad F_{223}=x\beta+a\,G_{kk},
\end{equation}
\begin{center}
\begin{tabular}{@{}l@{\quad}l@{\qquad}l@{\quad}l@{\qquad}l@{\quad}l@{}}
$\fsl_3$ & $(-\frac23,-\frac1{12})$ & $\fsl_4$ & $(-\frac38,-\frac3{64})$ & $\fsl_5$ & $(-\frac6{25},-\frac3{100})$\\
$\fsl_6$ & $(-\frac16,-\frac1{48})$ & $\fso_5=\fsp_4$ & $(-1,-\frac1{18})$ & $\fso_7$ & $(-\frac3{25},-\frac9{250})$\\
$\fsp_6$ & $(-\frac{21}{32},-\frac{15}{512})$ & $\fso_8$ & $(0,-\frac1{36})$ & $\fso_9$ & $(\frac3{49},-\frac{15}{686})$\\
$\fsp_8$ & $(-\frac{12}{25},-\frac9{500})$ & $\fso_{10}$ & $(\frac3{32},-\frac9{512})$ & $G_2$ & $(-\frac7{24},-\frac5{96})$\\
$F_4$ & $(\frac4{27},-\frac5{324})$ & $E_6$ & $(\frac5{24},-\frac1{96})$ & $E_7$ & $(\frac7{27},-\frac1{162})$\\
$E_8$ & $(\frac{22}{75},-\frac1{300})$ & & & &
\end{tabular}
\end{center}
(pairs $(x,a)$). For $\fsl_2$: $F_{133}=\frac72\beta$, $F_{223}=-3\beta$, $-3F_{133}+2F_{223}=-\frac{33}2\beta$. Hence, for $\g\in\mathcal L$, $[F_{133}]=-[F_{223}]$ and
\begin{equation}
O_2=\frac{\pi\cN}{48}F_{223}-\frac{\pi\cN}{160}\beta,\qquad [O_2]=\frac{\pi\cN}{48}[F_{223}]=\frac{\pi\cN}{48}\,a\,[G_{kk}]\neq0;\qquad \fsl_2:\ [O_2]=0 .
\end{equation}
For $\kappa={\rm Killing}/r$: $F_{133}+F_{223}=\frac{r^2}2\beta_\kappa$; values $\frac{r^2}{2}$ in trace forms: $\fso_5$ $\frac92$ ($\tr_5$), $\fsp_4$ $18$ ($\tr_4$), $\fso_6$ $8$ ($\tr_6$), $\fso_7$ $\frac{25}2$ ($\tr_7$), $\fsp_6$ $32$ ($\tr_6$), $\fso_8$ $18$ ($\tr_8$), $G_2$ $8$ ($\tr_7$), $\fsl_3,\fsl_4,\fsl_5$ $18,32,50$ ($\tr_N$). Under $\kappa\to c\kappa$: $W,F,\omega,O_2\to c^{-1}(\cdot)$, $\beta\to c\beta$.
\end{theorem}
\cert{C16, C18, C20}

\begin{proposition}[invariant dimensions]\label{prop-inv}
Weyl-formula dimensions of invariants:
{\small
\begin{longtable}{@{}lrrrrrrrr@{}}
\toprule
 & $S^2$ & $\Lambda^2$ & $S^3$ & $\Lambda^3$ & $T^3$ & $S^4$ & $\Lambda^4$ & $T^4$\\
\midrule
$\fsl_2$ & 1&0&0&1&1&1&0&3\\
$\fsl_3$ & 1&0&1&1&2&1&0&8\\
$\fsl_4=\fso_6$ & 1&0&1&1&2&2&0&9\\
$\fsl_5$ & 1&0&1&1&2&2&0&9\\
$\fso_5=\fsp_4$ & 1&0&0&1&1&2&0&6\\
$\fso_7$ & 1&0&0&1&1&2&0&6\\
$\fsp_6$ & 1&0&0&1&1&2&0&6\\
$\fso_8$ & 1&0&0&1&1&3&0&7\\
$G_2$ & 1&0&0&1&1&1&0&5\\
$F_4$ & 1&0&0&1&1&1&0&5\\
\bottomrule
\end{longtable}}
Further: $\dim(S^3\g^*)^\g=0$, $\dim(\Lambda^3\g^*)^\g=\dim(\g^{*\otimes3})^\g=1$ for $\fso_N$ ($5\le N\le40$, $N\ne6$), $\fsp_N$ ($N$ even, $4\le N\le40$), $G_2,F_4,E_6,E_7,E_8$; $\dim(S^3\g^*)^\g=\dim(\Lambda^3\g^*)^\g=1$, $\dim(\g^{*\otimes3})^\g=2$ for $\fsl_N$, $3\le N\le30$. Over the table: $[O_2]\neq0\iff\dim(\g^{\otimes4})^\g>3$; for $\fsl_2$: $C^4_{\rm inv}=\mathbb C\beta$.
\end{proposition}
\cert{C17, C19}

\begin{dstatement}[single computation]\label{ds-C4inv}
$\dim C^4_{\rm inv}=1$ ($\fsl_2$), $3$ ($\fsl_N$, $N=3,4,5$), $2$ ($\fso_5$, $\fso_7$, $\fsp_6$, $\fso_8$, $G_2$, $F_4$) (character computation); rank of $d$ on $C^4_{\rm inv}$: $0$ ($\fsl_2,\fso_5,G_2$), $1$ ($\fsl_3$); $H^4_{\rm inv}=\mathbb C[F_{223}]$ for $\fsl_3,\fso_5,G_2$, $H^4_{\rm inv}(\fsl_2)=0$.
\end{dstatement}

\subsubsection*{All simple Lie algebras}

Probe. For $X=\sum_\mu\mu(p)X^\mu\in C^4_{\rm inv}$, $\mu$ in the monomial basis $(u^2,v^2,w^2,uv,uw,vw)$, $u=[p_0p_1]$, $v=[p_0p_2]$, $w=[p_1p_2]$:
\begin{gather}
\ell_{kk}(X):=\langle X^{uv},kk_{(0123)}\rangle,\qquad \ell_{ff}(X):=\langle X^{uv},ff_{(0123)}\rangle,\\ kk_{(0123)}(x)=\kappa(x_0,x_1)\kappa(x_2,x_3),\quad ff_{(0123)}(x)=\kappa([x_0,x_1],[x_2,x_3]),\\
\varpi(X):=\ell_{kk}(X)\,\ell_{ff}(\beta)-\ell_{ff}(X)\,\ell_{kk}(\beta),
\end{gather}
$\langle\cdot,\cdot\rangle$: full contraction with $\kappa^{-1}$. $\varpi(\beta)=0$; by Lemma~\ref{lem-red}, $X$ exact $\Rightarrow$ $X\in\mathbb C\beta$ $\Rightarrow$ $\varpi(X)=0$. Transfer, for $X=\alt([p_0p_1][p_2p_3]W)$, $w^{kk}_\pi=\langle W,kk_\pi\rangle$, $w^{ff}_\pi=\langle W,ff_\pi\rangle$, $\pi\in((0123),(0213),(0312))$:
\begin{gather}
\ell_{kk}(X)=8w^{kk}_0-4w^{kk}_1-4w^{kk}_2,\qquad \ell_{ff}(X)=-4w^{ff}_1-4w^{ff}_2,\\
\ell_{kk}(\beta)=2d,\quad \ell_{ff}(\beta)=d,\quad \ell_{kk}(G_{kk})=8d(d-1),\quad \ell_{ff}(G_{kk})=8d\qquad(\text{Killing},\ d=\dim\g).
\end{gather}
Networks, $\Omega_{\rm ad}=\sum_x{\rm ad}(e_x)\otimes{\rm ad}(e^x)$ on $\g\otimes\g$, $P$ the exchange, $m_k=\tr(\Omega_{\rm ad}^k)$, $n_k=\tr(P\Omega_{\rm ad}^k)$:
\begin{gather}
W_{223}:\ w^{kk}=(-\tfrac d2,\tfrac d4,\tfrac d4),\ w^{ff}=(0,n_4,n_4);\qquad W_{133}:\ w^{kk}=(d,\tfrac d4,0),\ w^{ff}=(-\tfrac d8,-m_4,-n_4);\\
n_2=\tfrac d2\ \Big(\textstyle\sum_y{\rm ad}_y{\rm ad}_X{\rm ad}^y=\tfrac12{\rm ad}_X\Big),\qquad m_3=-\tfrac d4,\qquad m_4-n_4=\tfrac d8,\\
\ell_{kk}(F_{223})=-6d,\quad \ell_{ff}(F_{223})=-8n_4,\qquad \varpi(F_{223})=2d(8n_4-3d),\\ \varpi(F_{133})=-\varpi(F_{223}),\qquad \varpi(G_{kk})=8d^2(d-3).
\end{gather}
$n_4=\sum_{V\subset S^2\g}\dim V\,w_V^4-\sum_{V\subset\Lambda^2\g}\dim V\,w_V^4$, $w_V=(C_2(V)-2C_2(\g))/2$.

\begin{theorem}[all simple $\g$]\label{thm-allsimple}
$\g$ simple, $\kappa$ Killing.
(i) $8n_4-3d$: $\fso_N$: $-\frac{3N(N-1)(N-3)(N-4)(N+2)}{2(N-2)^3}$; $\fsp_N$: $-\frac{3N(N-2)(N+1)(N+3)(N+4)}{2(N+2)^3}$; $\fsl_N$: $-\frac{3(N^2-1)(N^2-4)}{N^2}$; $G_2$ $-\frac{385}{12}$, $F_4$ $-\frac{12740}{81}$, $E_6$ $-\frac{975}4$, $E_7$ $-\frac{34580}{81}$, $E_8$ $-\frac{12152}{15}$. Integer zeros: $\fso_3=\fsl_2$, $\fso_4$ (not simple), $\fsp_2=\fsl_2$, $\fsl_2$.
(ii) $\varpi(o_2)=\frac{\pi\nu}{48}\varpi(F_{223})=\frac{\pi\nu}{24}\,d\,(8n_4-3d)$; for $o_2$ built with $-3F_{133}+2F_{223}$, $-3F_{133}-2F_{223}$, $3F_{133}+2F_{223}$: $\varpi=5,\,1,\,-1$ times $\frac{\pi\nu}{240}\varpi(F_{223})$.
(iii) $o_2\notin d(C^3)$ for every simple $\g\neq\fsl_2$; for $\fsl_2$, $o_2=-\frac{33}2\frac{\pi\nu}{240}\beta$ is exact.
(iv) Vogel parameters (Mkrtchyan--Sergeev--Veselov 1105.0115v3 \cite{MSV}, Table~1, eqs.~(1), (8), (9), p.~4) $(v_1,v_2,v_3)$, $t_V=v_1+v_2+v_3$: $n_4=d(C_4-\frac18)$, $C_4=\frac{3t_Vt_2-t_3}{16t_V^3}$,
\begin{equation}
\begin{gathered}8n_4-3d=-\frac{3d(t_V+v_1)(t_V+v_2)(t_V+v_3)}{2t_V^3}=\frac{3d(d-3)v_1v_2v_3}{4t_V^3},\\ \varpi(F_{223})=\frac{3d\,\dim(X_2)\,v_1v_2v_3}{t_V^3},\quad \dim X_2=\tfrac{d(d-3)}2,\end{gathered}
\end{equation}
and on $\mathcal L$: $a=\frac{\varpi(F_{223})}{\varpi(G_{kk})}=\frac{3v_1v_2v_3}{16t_V^3}$, $x=-3-4(d-1)a$; $a(\fso_N)=-\frac{3(N-4)}{2(N-2)^3}$, $a(\fsp_N)=-\frac{3(N+4)}{2(N+2)^3}$, $a(\fsl_N)=-\frac3{4N^2}$.
\end{theorem}
Proof. (i) from the families $\kappa=\tr_V$:
\begin{equation*}
n_4^{\rm tr}(\fso_N)=-\tfrac{3N(N-1)(N-2)(N^2-14N+32)}{16},\quad n_4^{\rm tr}(\fsp_N)=\tfrac{3N(N+1)(N+2)(N^2+14N+32)}{16},\quad n_4^{\rm tr}(\fsl_N)=24N^2(N^2-1),
\end{equation*}
Killing $=(N-2),(N+2),2N$ times $\tr_V$, $n_4=n_4^{\rm tr}/r^4$; $n_4^{\rm tr}$ is a polynomial of degree $\le8$ in $N$ (traces over $V^{\otimes4}$ of permutations and at most four contractions) and agrees with the computed values at more than nine $N$. (ii) from the network values. (iii) $o_2$ is $G\times SL_2$-invariant; $o_2\in d(C^3)$ $\Rightarrow$ $o_2=d(R\eta)$, $R\eta\in[q_0q_1]^2\otimes(\Lambda^3\g^*)^\g$ (Lemma~\ref{lem-red}(a),(b)); $\varpi(d([q_0q_1]^2T))=0$ for $T\in(\Lambda^3\g^*)^\g$: for $T=f$ by $\varpi(\beta)=0$, for $T=cf+T'$, $\langle T',f\rangle=0$, the contractions reduce by $\kappa^{ij}f(i,j,\cdot)=0$, $ff\to C\kappa$ and $fff\to\frac C2f$ to multiples of $\langle T',f\rangle$ (this case is not needed where $\dim(\Lambda^3\g^*)^\g=1$, Proposition~\ref{prop-inv}); then (i), (ii) give $\varpi(o_2)\ne0$. $\fsl_2$ from Theorem~\ref{thm-simple}. (iv) $t_2=t_V^2-2s$, $t_3=t_V^3-3t_Vs+3p$, $(t_V+v_1)(t_V+v_2)(t_V+v_3)=2t_V^3+t_Vs+p$, $d\,p=p-2t_Vs-4t_V^3$ ($s,p$ the second and third elementary symmetric functions of $v_i$); $\varpi(F_{223})=0\iff\dim X_2=0$ ($d=3$) or $v_1v_2v_3=0$, the latter only for $\fso_4$. With Derived statement~\ref{ds-closed} (closedness), $[O_2]\neq0$ in $H^1(V)$ for every simple $\g\ne\fsl_2$ (Corollary~\ref{cor-O2}).
\cert{C18, C21, C22, C23}

%======================================================================
\section{\texorpdfstring{Orders $\hbar$ and $\hbar^2$}{Orders hbar and hbar2}}\label{sec-deform}
%======================================================================

\subsection{Quantizations considered}

\begin{definition}[setting]\label{def-setting}
A quantization through order $\hbar^n$: $S=\sum_{k=0}^n\hbar^kI_k$, $I_0=S_{\rm cl}$, $I_1=J_1$, $I_k\in\Vht$ of ghost number $0$ ($k\ge1$), regularized by the heat-kernel scheme of Section~\ref{sec-reg}, satisfying the quantum master equation through order $\hbar^n$ as $L\to0$. $U(2)$-equivariant: $S/\hbar$ invariant under $A\to\phi_g^*A$, $\hbar\to\det(g)^{-1}\hbar$, $g\in U(2)$, i.e.\ $I_k(\phi_g^*A)=\det(g)^{k-1}I_k(A)$.
\end{definition}
For $I_k\in\Vht$ with cochain of bidegree $(a,b)$: $I_k(\phi_l^*A)=l_1^{a-1}l_2^{b-1}I_k(A)$, so $U(2)$-equivariance $\iff$ $I_k$ of bidegree $(k,k)$ and $SL_2$-invariant (on polynomial cochains $SU(2)$- and $SL_2(\mathbb C)$-invariance coincide). $J_1$: bidegree $(1,1)$; $I_2$, $O_2$: $(2,2)$; $I_3$, $O_3$: $(3,3)$. By Theorem~\ref{thm-B} and Proposition~\ref{prop-eta}, cohomological statements about $\Vht$, $\Veta$ are statements about $\Cg$ with ghost number $g\leftrightarrow$ CE degree $g+3$; by Theorem~\ref{thm-C} $G$-invariance of the $I_k$ is no restriction on classes.

\subsection{\texorpdfstring{QME at orders $\hbar$, $\hbar^2$}{QME at orders hbar, hbar2}}\label{sec-qme}

S1 Sec.~9.3 (finite action deformations and their proof): $\mathrm{QME}_L(I+\delta)=\mathrm{QME}_L(I)+\mathbf Q_L\delta+\frac12(\delta,\delta)_L$, $\mathbf Q_L=Q+\hbar\Delta_L+(I[L],\cdot)_L$; $\mathbf Q_L=Q_0+\hbar Q_{(1)}+\dots$, $Q_0\to\{S_{\rm cl},\cdot\}$, $Q_{(1)}=\Delta_L+(I^{(1)}[L],\cdot)_L$.

\begin{proposition}[QME at orders $\hbar$, $\hbar^2$]\label{prop-qme}
With $S=S_{\rm cl}+\hbar J_1+\hbar^2I_2$, $J_1,I_2\in\Vht$:
\begin{equation}
\hbar^1:\ \{S_{\rm cl},J_1\}=0,\qquad \hbar^2:\ \{S_{\rm cl},I_2\}+\tfrac12\{J_1,J_1\}+O_2/\hbar^2=0 .\label{eq-qme2}
\end{equation}
\end{proposition}
$O_1=0$, $Q_{(1)}J_1=0$: these are $b_1=1$ graphs; a holomorphic derivative on a vertex inserts a polynomial in the currents $y_e$ (Section~\ref{sec-o3-signs}), a $\tau$-independent $0$-form; the $\tau$-factor of the reduced kernel is $\rho_1=0$, equivalently the reflection argument of Section~\ref{sec-reg} with $s$ commuting with $\partial_z$. The factor $\frac12$ relative to $O_2$: $O_2$ is the $\hbar^2$ part of $\mathrm{QME}_L(I[L])$ as $L\to0$ (Section~\ref{sec-reg}).
\cert{C34}

\subsection{Taylor cochains and bracket}

\begin{gather}
\{a,b\}(x_1..x_{m+n-2})=\sum_{s\in{\rm Sh}(m-1,n-1)}\sgn(s)\,a(x_{s(1)}..x_{s(m-1)},X_e)\,\kappa^{ee'}\,b(X_{e'},x_{s(m)}..),\\
\{a,b\}=\frac{\alt(\gamma_0)}{(m-1)!(n-1)!},\qquad
\gamma_0=a(x_1..x_{m-1},X_e)\,\kappa^{ee'}\,b(X_{e'},y_1..y_{n-1}),\\ \text{contracted momenta }q,-q\text{ fixed by conservation};\\
\fgl_N:\ \tr(AX_e)\tr(BX^e)=\tr(AB),\ \tr(X_eAX^eB)=\tr A\tr B;\\
\fsl_N:\ \text{additionally }-\tfrac1N\tr A\tr B,\ -\tfrac1N\tr(AB);\\
s_0(x_1,x_2,x_3)=\kappa(X_1,[X_2,X_3])=\tr(123)-\tr(132)\qquad(\text{cubic Taylor cochain of }S_{\rm cl}).
\end{gather}

\begin{proposition}[cochain form]\label{prop-cochain}
$\{s_0,w\}=-dw$; $\{a,b\}=\{b,a\}$ on arity $(3,3)$; $d\{a,b\}=\{da,b\}-\{a,db\}$ ($a$ of arity 3); $d^2=0$. With $\{F_a,F_b\}=\epsilon\,F_{\{a,b\}}$ (Proposition~\ref{prop-tree}) and $j,i_2$ the cochains of $J_1,I_2$, \eqref{eq-qme2} reads
\begin{equation}
\epsilon\big(-d\,i_2+\tfrac12\{j,j\}\big)+o_2=0,\qquad \big[\tfrac12\{j,j\}\big]=-\epsilon\,[o_2]\ \ \text{in }H^4_{(2,2)} .\label{eq-cochain}
\end{equation}
\end{proposition}
$\{s_0,w\}$: the $(2,k-1)$ shuffle putting $(x_i,x_j)$ in front has sign $(-1)^{i+j-1}$ and $s_0(x_i,x_j,X_e)\kappa^{ee'}=\kappa([x_i,x_j],X_e)\kappa^{ee'}$. Dictionary factor $\frac1{r!}\{a,b\}(c..c)=\gamma_0(c..c)/((m-1)!(n-1)!)$.
\cert{C35}

\begin{proposition}[tree graph and the bracket dictionary]\label{prop-tree}
Tree graph with one heat edge $e=(v_0,v_1)$, cochain $a$ at $v_0$ (externals $x_1,x_2$), $b$ at the root $v_1$ (externals $x_3,x_4$), target $((e,0),(e,1),x_1,x_2,x_3,x_4)$, $G_{\rm tree}=\int_X\Omega\wedge K_\varepsilon=1$, $|\Aut|=4$, $y_e=p_1+p_2$, Wick rule of Proposition~\ref{prop-wick}:
\begin{equation}
\tfrac14\alt\big(\sigma\,a(L_0)\,\kappa^{-1}\,b(L_1)\big)=\{a,b\}\quad\Longrightarrow\quad \{F_a,F_b\}=U_{\rm tree}\,F_{\{a,b\}},\qquad \epsilon=U_{\rm tree},
\end{equation}
for $(a,b)=(D,\uf),(D,D),(s_0,\uf),(S_1,\uf)$ (notation of Sections~\ref{sec-H3}, \ref{sec-o3}), all $36$ pairs of local orders.
\end{proposition}
\cert{C39}

\begin{dstatement}[$U_{\rm tree}=U$]\label{ds-epsU}
Connected graph, $m$ cubic ghost-$0$ vertices, $|I|$ edges (one heat edge): $\prod_{v=1}^m\int_v\omega_v=(-1)^{\binom m2}\int_1..\int_m\omega_1..\omega_m$; Koszul reordering of the $A$'s to the target: $\sigma_\Gamma$; pairs $\to s_P\kappa^{-1}P_e$ (even), heat pair $\to s_K\kappa^{-1}K_\varepsilon$ (odd), $s_K=s_hs_P$ ($s_h$: sign of $DP=K_\varepsilon-K_L$ relative to the $\Delta_\varepsilon$ contraction), product $s_hs_P^{|I|}$; $\int_m$ past $\Omega_1..\Omega_{m-1}(\prod P\wedge K)$: $(-1)^{5(m-1)}$; $\int_1..\int_{m-1}=(-1)^{\binom{m-1}2}\int_{X^{m-1}}$;
\begin{equation}
\tbinom m2+\tbinom{m-1}2+5(m-1)=(m-1)(m+4)\equiv0\pmod2\ \Rightarrow\ w_\Gamma=U_0\,s_P^{|I|}\hbar^{b_1}\sigma_\Gamma\int_X\Omega\,[G_\Gamma C_\Gamma A^E],
\end{equation}
$U_0$ graph-independent; tree $(m,|I|)=(2,1)$, thetas $(6,7)$: $U_{\rm tree}=U_0s_P=U_0s_P^7=U$, hence $\epsilon U=+1$. Checked: the exponent identity and the parities of $|I|$; the signs $s_h$, $s_P$ and the placement of the external block are not recomputed.
\end{dstatement}

\subsection{\texorpdfstring{$H^3$ at weight $(1,1)$ and at all bidegrees}{H3 at weight (1,1) and at all bidegrees}}\label{sec-H3}

\begin{theorem}[closed order-$\hbar$ terms]\label{thm-H3}
For $\fgl_N,\fsl_N$ ($\kappa=\tr$), at bidegree $(1,1)$ and at total weight $2$:
\begin{equation}
H^3_{(1,1)}=\{[p_1p_2]\,T:\ T\in(\mathrm{Sym}^3\g^*)^\g\},\qquad H^3_{(2,0)}=H^3_{(0,2)}=0,\qquad \alt([p_1p_2]T)=6[p_1p_2]\,\mathrm{Sym}(T).
\end{equation}
$\fgl_N$ ($N\ge3$) basis: $D=[p_1p_2]\tr(1\{23\})$, $S_1=[p_1p_2](\tr1\tr23+\tr2\tr13+\tr3\tr12)$, $S_0=[p_1p_2]\tr1\tr2\tr3$; $\fsl_N$ ($N\ge3$): $D$. $(\dim C^3_{(1,1)},\dim H^3_{(1,1)})$: $\fgl_1$ $(1,1)$, $\fgl_2$ $(4,2)$ [$D=S_1-S_0$], $\fgl_{3,4,5}$ $(5,3)$, $\fsl_2$ $(1,0)$, $\fsl_{3,4,5}$ $(2,1)$. $E_f=\alt(q_0^1q_0^2f)=2(\sum_ip^1_ip^2_i)\tr(1[23])$, $E_{\rm mix}=\alt((p^1_1p^2_2+p^2_1p^1_2)\tr1\tr23)$ are not closed.
Simple $\g$, $G\times SL_2$-invariant weight $(1,1)$: every such 3-cochain is $[q_0q_1]T$, $T\in(S^3\g^*)^\g$; all are closed, none nonzero is exact.
\end{theorem}
$C^2=0$ at bidegrees $(2,0),(1,1),(0,2)$ (the momentum monomial of degree $2$ in one variable is even under $1\leftrightarrow2$, the invariant 2-tensors are symmetric), so $H^3=Z^3$ there; $[p_{s(1)}p_{s(2)}]=\sgn(s)[p_1p_2]$ modulo $\sum p$.
\cert{C36}

\begin{proposition}[all bidegrees, $\dim(\g^{*\otimes3})^\g=1$]\label{prop-H3all}
$\g$ simple with $(\g^{*\otimes3})^\g=\mathbb Cf$: $C^2_{(a,b)}={\rm span}\,\alt(q_0^m\kappa)$ (nonzero iff $a+b$ odd), $C^3_{(a,b)}=({\rm Sym}\,m)f$, $d(({\rm Sym}\,m)f)=c_1(p)T_1+c_3(p)T_3$, $T_1=ff_{(0123)}$, $T_3=ff_{(0312)}$, $ff_{(0213)}=T_1+T_3$; $T_1,T_3$ linearly independent ($xT_1+yT_3=0$ with the exchange gives $T_1=-2T_3$, contraction with $\kappa^{ab}\kappa^{ce}$ gives $0=2d$). The table (G-invariant cochains, $0\le a,b\le2$) is independent of $\g$:
\begin{center}
\begin{tabular}{@{}lccccccccc@{}}
\toprule
$(a,b)$ & $(0,0)$ & $(0,1)$ & $(1,0)$ & $(0,2)$ & $(2,0)$ & $(1,1)$ & $(1,2)$ & $(2,1)$ & $(2,2)$\\
\midrule
$\dim C^2$ & 0&1&1&0&0&0&1&1&0\\
$\dim C^3$ & 1&0&0&1&1&1&1&1&2\\
${\rm rank}\,d_2$ & 0&0&0&0&0&0&1&1&0\\
${\rm rank}\,d_3$ & 0&0&0&1&1&1&0&0&2\\
$\dim H^3$ & 1 ($s_0$)&0&0&0&0&0&0&0&0\\
\bottomrule
\end{tabular}
\end{center}
At $(1,1)$: $C^3=\mathbb CE_f$,
\begin{equation}
\begin{gathered}dE_f=2\sum_{i<j}(-1)^{i+j}(p_i^1p_j^2+p_i^2p_j^1)\,\kappa([X_i,X_j],[X_k,X_l])=2\big[B(p_0-p_3,p_2-p_1)T_1+B(p_0-p_1,p_2-p_3)T_3\big]\neq0,\\ B(p,q)=p^1q^2+p^2q^1 .\end{gathered}
\end{equation}
Contrast $\fsl_3,\fsl_4$: $(1,1)$: $C^3$ $2$, ${\rm rank}\,d_3$ $1$, $H^3$ $1$; $(2,2)$: $C^3$ $3$, ${\rm rank}\,d_3$ $3$, $H^3$ $0$.
\end{proposition}
\cert{C24, C25}

\subsection{Moyal term}

\begin{proposition}[first orders of the Moyal product]\label{prop-moyal}
$f*g=fg+\frac e2\epsilon^{ij}\partial_if\partial_jg+O(e^2)$, $e_p*e_q=e^{\frac e2[pq]}e_{p+q}$; the $\fgl_N$ cubic cochain of $\frac16\int\Omega\,\kappa(A,[A,A]_*)$:
\begin{gather}
s_e=\tfrac13\alt\big(\tr(123)e^{\frac e2[p_1p_2]}\big)=\tr(123)e^{\frac e2[p_1p_2]}-\tr(132)e^{-\frac e2[p_1p_2]}=s_0+e\,j_M+e^2i_M+O(e^3),\\
j_M=\tfrac12[p_1p_2]\tr(1\{23\})=\tfrac12D,\qquad i_M=\tfrac18[p_1p_2]^2\tr(1[23]),\\
dj_M=0,\qquad \tfrac12\{j_M,j_M\}=d\,i_M,\qquad\{j_M,j_M\}\ne0 .
\end{gather}
$\pi X=X-\tr X/N$: $\pi^*D_{\fsl}=D-\frac2NS_1+\frac4{N^2}S_0$, $M_1:=S_1-\frac3NS_0$, $j_M=\frac12\pi^*D_{\fsl}+\frac1NM_1+\frac1{N^2}S_0$.
\end{proposition}
Normalization: Costello 1610.04144 \cite{C16} Sec.~5, $f*_eg=fg+e\,\frac12\epsilon^{ij}\partial_if\partial_jg+\dots$, coefficient of $e^n$: $(2^nn!)^{-1}\epsilon\cdots\epsilon$; Sec.~9.1: $[z_1,z_2]=e$. This is the convention above.
\cert{C37}

\subsection{\texorpdfstring{The quadratic map $q[J]=[\frac12\{J,J\}]$}{The quadratic map q}}

\begin{theorem}[brackets of $H^3$ classes]\label{thm-quad}
Formal $\fgl_N$, exact cochain identities ($\beta=d\,\uf$, $\uf:=[p_1p_2]^2\tr(1[23])$, $F_{223}=F_0+NF_1$):
\begin{gather}
\tfrac12\{D,D\}=\tfrac12\beta,\quad\{D,S_1\}=-\tfrac1{12}F_1,\quad\{D,S_0\}=\tfrac1{24}F_0+\beta,\\
\tfrac12\{S_1,S_1\}=\tfrac1{48}F_0-\tfrac N{48}F_1+\tfrac12\beta,\quad\{S_1,S_0\}=\tfrac N{48}F_0+\tfrac N2\beta,\quad\{S_0,S_0\}=0;\\
J=aD+bS_1+cS_0:\quad \tfrac12\{j,j\}=\lambda F_{223}+\mu F_1+k_\beta\beta,\\
\lambda=\frac{2ac+b^2+Nbc}{48},\qquad \mu=-\frac{(2a+Nb)(2b+Nc)}{48},\qquad k_\beta=\frac{a^2+2ac+b^2+Nbc}2;\\
J=\alpha\,\pi^*D_{\fsl}+\chi M_1+\gamma S_0\ (a=\alpha,\ b=\chi-\tfrac{2\alpha}N,\ c=\gamma-\tfrac{3\chi}N+\tfrac{4\alpha}{N^2}):\\ \lambda=\frac{4\alpha^2-2N^2\chi^2+N^3\chi\gamma}{48N^2},\qquad \mu=-\frac{N\chi(N\gamma-\chi)}{48}.
\end{gather}
Classes in the basis $\{[F_{223}],[F_1],[z_8],[z_9]\}$: $q[J]=\lambda[F_{223}]+\mu[F_1]$. Formal $\fsl_N$: $\{D,D\}=\frac1{6N}[F_1]$, $q[aD]=\frac{a^2}{12N^2}[F_{223}]$, with
\begin{equation}
\tfrac12\{D,D\}_{\fsl}=\frac{F_{223}}{12N^2}+d\eta_N,\qquad \eta_N=\frac{N^2+4}{2N^2}\,\uf,\qquad \big[\alt([p_1p_2][p_3p_4]\tr12\tr34)\big]=-\frac{[F_{223}]}{6N}.
\end{equation}
Fixed $N$ ($x_0,x_1,x_2$ the coefficients of $D,S_1,S_0$): $\fsl_3,\fsl_4,\fsl_5$: $q=\frac{x_0^2}{108},\frac{x_0^2}{192},\frac{x_0^2}{300}[F_{223}]$; $\fgl_N$, $N=3,4,5$: $\frac{2x_0x_2+x_1^2+Nx_1x_2}{48}[F_{223}]-\frac{(2x_0+Nx_1)(2x_1+Nx_2)}{48}[F_1]$; $\fgl_2$ ($J=x_0D+x_1S_1$): $q=-\frac{x_1(x_0+x_1)}{12}[F_1]$, $[F_1]\ne0$, $[F_{223}]=0$; $\fgl_1$: $q=0$.
\end{theorem}
$\lambda$ refers to the complement $\{[F_1],[z_8],[z_9]\}$ of $[F_{223}]$. Family $s\,j_M+t\,\pi^*D_{\fsl}$: $(\alpha,\chi,\gamma)=(\frac s2+t,\frac sN,\frac s{N^2})$, $2a+Nb=s$, $2b+Nc=0$, $\lambda=\frac{t(s+t)}{12N^2}$.
\cert{C38, C46}

\subsection{\texorpdfstring{Cancellation of $[O_2]$ and no-go}{Cancellation of [O2] and no-go}}

$c_O:=[o_2]/[F_{223}]=\frac{\pi\nu}{48}=-\frac{U}{3072\pi^4}$ (Theorem~\ref{thm-class-slN}). By \eqref{eq-cochain} the order-$\hbar^2$ equation is solvable iff $q[j]=-\epsilon c_O[F_{223}]$.

\begin{theorem}[order $\hbar^2$, $\fsl_N$ and $\fgl_N$, $\kappa=\tr$]\label{thm-Q1}
Setting of Definition~\ref{def-setting}, $U(2)$-equivariant.
(i) $\fsl_N$, $N\ge3$: solvable iff $j=a\,D+(\text{exact})$, $j=a[p_1p_2]d_{abc}$, $d_{abc}=\tr(T_a\{T_b,T_c\})$, $J_1=a\,\frac13\int dz_1dz_2\,d_{abc}\,\partial_{z_1}A^a\,\partial_{z_2}A^b\,A^c$, with $\frac{a^2}{12N^2}=-\epsilon c_O$, i.e.\ $a^2=\epsilon U\frac{N^2}{256\pi^4}$.
(ii) $\fgl_N$, $N\ge3$: solvable iff $j=\alpha\,\pi^*D_{\fsl}+\chi M_1+\gamma S_0+(\text{exact})$ with $\chi(N\gamma-\chi)=0$ and $\frac{\alpha^2}{12N^2}-\frac{\chi^2}{48}=-\epsilon c_O$; examples $j=a\,\pi^*D_{\fsl}$ ($\frac{a^2}{12N^2}=-\epsilon c_O$), $j=b(S_1-\frac2NS_0)$ ($-\frac{b^2}{48}=-\epsilon c_O$, $b^2=-\frac{\epsilon U}{64\pi^4}$), $j=s\,j_M+t\,\pi^*D_{\fsl}$ ($\frac{t(t+s)}{12N^2}=-\epsilon c_O$). $j=s\,j_M$ has $q=0$ and does not solve.
(iii) $\fgl_1,\fgl_2,\fsl_2$: $[O_2]=0$; $J_1=0$ solves \eqref{eq-qme2}; all solutions: $\fgl_1$ $x_0D$; $\fgl_2$ $x_0D+x_1S_1$, $x_1(x_0+x_1)=0$; $\fsl_2$ $j$ exact.
\end{theorem}
Over $\mathbb C$; over $\mathbb R$ the sign of $\epsilon U$ decides (Derived statement~\ref{ds-epsU}: $\epsilon U=+1$).
\cert{C38}

$I_2$ for $\fsl_N$, $j=aD$: Section~\ref{sec-o3-I2}.

\begin{theorem}[no-go for simple $\g$]\label{thm-nogo}
$\g$ simple, $\g\neq\fsl_2$, $\dim(\g^{*\otimes3})^\g=1$ (Proposition~\ref{prop-inv}: $\fso_N$, $5\le N\le40$, $N\ne6$; $\fsp_N$, $N$ even, $4\le N\le40$; $G_2,F_4,E_6,E_7,E_8$). In the setting of Definition~\ref{def-setting} without equivariance: for every closed $J_1\in\Vht$ (any bidegrees, not necessarily $G$-invariant) the order-$\hbar^2$ equation has no solution $I_2\in\Vht$. In particular there is no $U(2)$-equivariant quantization through order $\hbar^2$.
\end{theorem}
Proof. $J_1=\sum_{(a,b)}J^{(a,b)}$, closed; by Theorem~\ref{thm-C} and $\{dK,J\}=\pm d\{K,J\}$ for closed $J$ (graded Jacobi with $\{s_0,\cdot\}=-d$) reduce to $G$-invariant $J$. Proposition~\ref{prop-H3all}: $J^{(1,0)}=J^{(0,1)}=0$ ($C^3=0$), $J^{(2,0)}=J^{(0,2)}=J^{(1,1)}=J^{(2,2)}=0$ ($Z^3=0$), $J^{(0,0)}=cs_0$. Bidegree is additive under the bracket; the weight-$(2,2)$ part
\begin{multline}
\tfrac12\!\!\sum_{(a,b)+(a',b')=(2,2)}\!\!\{J^{(a,b)},J^{(a',b')}\}=\{J^{(0,0)},J^{(2,2)}\}+\{J^{(1,0)},J^{(1,2)}\}+\{J^{(0,1)},J^{(2,1)}\}\\+\{J^{(2,0)},J^{(0,2)}\}+\tfrac12\{J^{(1,1)},J^{(1,1)}\}=0
\end{multline}
as a cochain; \eqref{eq-cochain} at weight $(2,2)$ becomes $o_2\in d(C^3)$, contradicting Theorem~\ref{thm-allsimple}(iii). $O_1=0$, $Q_{(1)}J_1=0$: Proposition~\ref{prop-qme}.
\cert{C24, C25, C18}

For $J_1\in\Veta$ not in $\Vht$ the statement uses Derived statement~\ref{ds-homotopy}.

%======================================================================
\section{\texorpdfstring{Order $\hbar^3$}{Order hbar3}}\label{sec-o3}
%======================================================================

Notation: $f=\tr(1[23])$, $d=\tr(1\{23\})$, $s_1=\tr1\tr23+\tr2\tr13+\tr3\tr12$, $s_0=\tr1\tr2\tr3$, $D=[p_1p_2]d$, $S_1=[p_1p_2]s_1$, $S_0=[p_1p_2]s_0$, $\uf=[p_1p_2]^2f$, $\Dthree=[p_1p_2]^3d$, $S_{1,3}=[p_1p_2]^3s_1$, $S_{0,3}=[p_1p_2]^3s_0$; $\pi\nu=-U/(64\pi^4)$; $b_{ij}=[p_ip_j]$. $\mathsf C:X\mapsto-X^{\mathsf T}$ (automorphism of $\fgl_N$, $\fsl_N$ preserving $\kappa=\tr$), on $k$-cochains $\tau\mapsto(-1)^k\tau^{\rm rev}$; $\mathsf C$ commutes with $d$ and the bracket.

\subsection{Equation and vanishing graphs}

\begin{proposition}[order-$\hbar^3$ equation]\label{prop-o3eq}
$S_{\rm int}=I+\hbar J_1+\hbar^2I_2+\hbar^3I_3$, $I_k\in\Vht$ ($QI_k=\int\Omega\,D(P_kW_k)=0$, $D\Omega=0$, $[D,\partial_z]=0$). A connected graph with vertices of orders $k(v)$ ($I$: $0$, $J_1$: $1$, $I_2$: $2$) carries $\hbar^{b_1+\sum_vk(v)}$. Order $\hbar^3$:
\begin{equation}
\{S_{\rm cl},I_3\}+\{J_1,I_2\}+\Theta(J_1)=0,\qquad \Theta(J_1)=\sum_{G\in\{\theta(1,3,3),\theta(2,2,3)\}}|\Aut G|^{-1}\sum_{v\in V(G)}w_{G,v},\label{eq-o3}
\end{equation}
$w_{G,v}$: the graph $G$ with $J_1$ at $v$ and $I$ elsewhere. In cochains: $\epsilon(-d\,i_3+\{j,i_2\})+\theta=0$, $\theta$ the cochain of $\Theta(J_1)/\hbar^3$. The remaining graphs vanish: $b_1=1$ with two $J_1$ or one $I_2$ vertex, $b_1=3$ with $I$ vertices only (odd $b_1$, Section~\ref{sec-reg}), self-loops.
\end{proposition}
Tree part: $\frac12\sum_{k+l=3}\{I_k,I_l\}=\{I,I_3\}+\{J_1,I_2\}$. Orbit--stabilizer: $\sum_{\Gamma'}|\Aut\Gamma'|^{-1}=\sum_G|\Aut G|^{-1}\sum_v$. The $J_1$ vertex has ghost number $0$ and $\partial_z$ is even, so $\Theta(J_1)$ carries the sign $U$ of $O_2$.
\cert{C34}

\subsection{\texorpdfstring{$J_1$ vertex and derivative rule}{J1 vertex and derivative rule}}\label{sec-o3-signs}

\begin{proposition}[vertex and derivative rules]\label{prop-o3rules}
(i) $F_D=\frac13\int\Omega\,d_{abc}\partial_{z_1}A^a\partial_{z_2}A^bA^c$ has Taylor cochain $\alt(\frac13p_1^1p_2^2d)=D$. Summing the $3!$ assignments to the local order $(h_1,h_2,h_3)$ with Koszul signs: $\frac13d_{h_1h_2h_3}\sum_s\sgn(s)q^1_{h_{s1}}q^2_{h_{s2}}=[q_{h_1}q_{h_2}]d_{h_1h_2h_3}$ on $\sum q=0$; the $J_1$ vertex factor is $a[q_{h_1}q_{h_2}]d_{h_1h_2h_3}$, $q$ the momentum into the vertex: $q_{x_v}=p_v$, $q_{(e,0)}=-y_e$, $q_{(e,1)}=+y_e$, and the Wick sign is $\sigma_G$ of the same graph with $f$ at that vertex.
(ii) $\partial/\partial z^i_{e(0)}\Theta_e=-y^i_e\Theta_e$, $\partial/\partial z^i_{e(1)}\Theta_e=+y^i_e\Theta_e$ ($dy_e$, $d\eta_e$ independent of $z$); by \eqref{eq-holGauss} $\zb\to4\Lambda^{-1}\lambda$ in polynomials, so $y_e$ becomes the Kirchhoff current of $B^{\mathsf T}y=\lambda$, $ty=B\varphi$: factor $1$, sign $-1$ at the tail, $+1$ at the head. Conservation: $\sum_hq_h=\lambda_v-\sum_{e(0)=v}y_e+\sum_{e(1)=v}y_e=0$.
\end{proposition}
\cert{C39, C40}

\begin{gather}
w_{G,v}=U\hbar^3\sigma_G\int_X\Omega\,\big[G^{(v)}_GC^{(v)}_GA^4\big],\qquad G^{(v)}_G=-\frac{K^{(v)}_G}{64\pi^5},\qquad \sigma_GC^{(v)}_G=a\,[q_{h_1}q_{h_2}]\,W^{(v,d)}_G,\\
W^{(v,T)}_G=\sgn(L\to L')\prod_{u\neq v}f(L'_u)\,T(L'_v)\prod_e\kappa^{-1}_{(e,0)(e,1)},\qquad K^{(v)}_G(\lambda;0)=K_{G,\htk}(\lambda;0)\big[\text{insertion }[q_{h_1}q_{h_2}]\big],\\
\theta=a\,\pi\nu\,\theta'(d),\qquad \theta'(T):=\sum_G|\Aut G|^{-1}\sum_{v\in V(G)}\alt\big(K^{(v)}_G/\pi\,W^{(v,T)}_G\big),\quad T\in\{d,s_1,s_0\}.
\end{gather}

\subsection{\texorpdfstring{$I_2$ for $\fsl_N$}{I2 for sl(N)}}\label{sec-o3-I2}

\begin{proposition}[$I_2$, $\fsl_N$, $N\ge3$]\label{prop-I2}
$j=aD$, $a^2=-12N^2\epsilon c_O=-\epsilon\pi\nu N^2/4$:
\begin{gather}
\tfrac12\{D,D\}=\frac{F_{223}}{12N^2}+d\eta_N,\qquad F_{133}+F_{223}=2N^2\,d\uf,\\ o_2-c_OF_{223}=-\frac{\pi\nu}{80}(F_{133}+F_{223})=d\zeta,\quad \zeta=-\frac{\pi\nu N^2}{40}\uf,\\
o'-\tfrac{N^2}8\{D,D\}=d\iota,\qquad \iota=-\frac{3N^2+10}{20}\,\uf,\\
i_2=a^2\eta_N+\epsilon\zeta=\epsilon\pi\nu\,\iota=\epsilon U\,\frac{3N^2+10}{1280\pi^4}\,\uf,\\ I_2=\epsilon U\,\frac{3N^2+10}{1280\pi^4}\,\frac13\int dz_1dz_2\,f_{abc}\big(\partial_{z_1}^2A^a\partial_{z_2}^2A^b-\partial_{z_1}\partial_{z_2}A^a\,\partial_{z_1}\partial_{z_2}A^b\big)A^c .
\end{gather}
$\dim C^2_{(2,2)}=0$, $\dim C^3_{(2,2)}=3$, ${\rm rank}\,d_3=3$: $i_2$ is unique ($G$-invariant complex; in general up to $Q$-exact terms, Theorem~\ref{thm-C}). $\mathsf C(i_2)=i_2$.
\end{proposition}
Chain: $d\,i_2=\frac12\{j,j\}+\epsilon o_2=a^2(\frac{F_{223}}{12N^2}+d\eta_N)+\epsilon o_2=a^2d\eta_N+\epsilon(o_2-c_OF_{223})$; $a^2\eta_N=\epsilon U\frac{N^2+4}{512\pi^4}\uf$, $\epsilon\zeta=\epsilon U\frac{N^2}{2560\pi^4}\uf$, sum $\epsilon U\frac{3N^2+10}{1280\pi^4}\uf$. $\alt\big(\frac13((p_1^1)^2(p_2^2)^2-p_1^1p_1^2p_2^1p_2^2)f\big)=\uf$.
\cert{C43}

\subsection{\texorpdfstring{Kernels with a $J_1$ vertex}{Kernels with a J1 vertex}}

\begin{proposition}[$J_1$-vertex kernels]\label{prop-o3kernels}
$K^{(v)}_G(\lambda;0)$ in the engine labelling, local order $L'$, $b_{ij}=[p_ip_j]$, $p_4=-p_1-p_2-p_3$, $[p_3p_4]=b_{13}+b_{23}$; $(h_1,h_2)$ the first two entries of $L'_v$:
{\small
\begin{longtable}{@{}lll@{}}
\toprule
$v$ & $L'_v$ & $K^{(v)}/\pi$\\
\midrule
\multicolumn{3}{@{}l}{$\theta(1,3,3)$, $K/\pi=-b_{12}(b_{13}+b_{23})/20$}\\
0 & $(e_0,0),(e_1,0),(e_4,0)$ & $b_{12}(b_{13}+b_{23})(b_{12}-2b_{13}-b_{23})/252$\\
1 & $x_1,(e_1,1),(e_2,0)$ & $b_{12}(b_{13}+b_{23})(7b_{12}-b_{13})/420$\\
2 & $x_2,(e_2,1),(e_3,0)$ & $b_{12}(b_{13}+b_{23})(8b_{12}-b_{23})/420$\\
3 & $x_3,(e_4,1),(e_5,0)$ & $b_{12}(b_{13}+b_{23})(8b_{13}+7b_{23})/420$\\
4 & $x_4,(e_5,1),(e_6,0)$ & $-b_{12}(b_{13}+b_{23})(b_{12}-7b_{13}-8b_{23})/420$\\
5 & $(e_0,1),(e_3,1),(e_6,1)$ & $-b_{12}(b_{13}+b_{23})(2b_{12}-b_{13}-2b_{23})/252$\\
\multicolumn{3}{@{}l}{$\theta(2,2,3)$, $K/\pi=b_{12}(b_{13}+b_{23})/30$}\\
0 & $(e_0,0),(e_2,0),(e_4,0)$ & $b_{12}(b_{13}+b_{23})(5b_{12}-2b_{13}+2b_{23})/840$\\
1 & $x_1,(e_0,1),(e_1,0)$ & $b_{12}(b_{13}+b_{23})(b_{12}+2b_{13})/420$\\
2 & $x_2,(e_2,1),(e_3,0)$ & $-b_{12}(b_{13}+b_{23})(b_{12}-2b_{23})/420$\\
3 & $x_3,(e_4,1),(e_5,0)$ & $-b_{12}(b_{13}+b_{23})^2/84$\\
4 & $x_4,(e_5,1),(e_6,0)$ & $-b_{12}(b_{13}+b_{23})^2/84$\\
5 & $(e_1,1),(e_3,1),(e_6,1)$ & $b_{12}(b_{13}+b_{23})(9b_{12}+2b_{13}-2b_{23})/840$\\
\bottomrule
\end{longtable}}
Pair independence: $[q_{h_1}q_{h_2}]=[q_{h_2}q_{h_3}]=[q_{h_3}q_{h_1}]$ modulo $\sum p=0$ in all $12$ cases. In another local order the kernels change by the uniform sign $\sgn(L'\to L'')$, compensated by $W^{(v,T)}$.
\end{proposition}
\cert{C41}

\subsection{\texorpdfstring{Colour factors and $\Theta(J_1)$}{Colour factors and Theta(J1)}}

\begin{proposition}[colour factors and $\theta'$]\label{prop-theta}
$\sgn(L\to L')=-1$ for both graphs. $\fsl_N$: $W^{(v,d)}_{133}=N^2(\tr(1342)-\tr(1243))$ ($v=0,1,2$), $N^2(\tr(1243)-\tr(1342))$ ($v=3,4,5$); $W^{(v,d)}_{223}=0$ ($v=0..5$). $\fgl_N$: additional terms containing the trace of a single $X$. $\theta'(s_1)=\theta'(s_0)=0$. By pieces, $|\Aut G|^{-1}\alt(K^{(v)}/\pi\,W^{(v,d)})=d(x\Dthree+yS_{1,3})$:
\begin{gather}
\theta(1,3,3):\ v=0,5:\ (x,y)=(-\tfrac{N^2}{504},\tfrac N{252});\quad v=1,2,3,4:\ (-\tfrac{N^2}{315},\tfrac{13N}{2520});\quad\text{total }(-\tfrac{N^2}{60},\tfrac N{35}),\notag\\
\theta(2,2,3):\ v=1,2:\ (0,\tfrac N{420});\quad v=0,3,4,5:\ (0,0);\quad\text{total }(0,\tfrac N{210}),\notag\\
\theta'(d)=d\Big(-\frac{N^2}{60}\Dthree+\frac N{30}S_{1,3}\Big)\ (\fgl_N),\qquad \theta'(d)|_{\fsl_N}=-\frac{N^2}{60}\,d\Dthree,\\
d\Dthree(x_0..x_3)=\sum_{i<j}(-1)^{i+j}[p_kp_l]^3\tr\big([X_i,X_j]\{X_k,X_l\}\big),\qquad\{k<l\}=\{0,1,2,3\}\setminus\{i,j\},\notag\\
\Theta(J_1)|_{\fsl_N}=-\frac{a\pi\nu N^2}{60}\,d\Dthree=\frac{aUN^2}{3840\pi^4}\sum_{i<j}(-1)^{i+j}[p_kp_l]^3\tr\big([X_i,X_j]\{X_k,X_l\}\big).\notag
\end{gather}
The piece values are constant on $\Aut$-orbits ($\{0,5\},\{1,2,3,4\}$ for $\theta(1,3,3)$; $\{0,5\},\{1,2\},\{3,4\}$ for $\theta(2,2,3)$). $\mathsf C(\theta')=-\theta'$.
\end{proposition}
Internal $\fsl_N$ propagators: the $\mathfrak u(1)$ part of $\kappa^{-1}_{\fsl}=\kappa^{-1}_{\fgl}-\frac1N\mathbf1\otimes\mathbf1$ meets an $f$ at the other end, $f(\mathbf1,\cdot,\cdot)=0$.
\cert{C42, C43}

\subsection{Closedness}

\begin{proposition}[closedness of $o_3$]\label{prop-o3closed}
$o_3:=\epsilon\{j,i_2\}+\theta$. $d\{j,i_2\}=\{dj,i_2\}-\{j,d\,i_2\}$, $dj=0$, $d\,i_2=\frac12\{j,j\}+\epsilon o_2$ $\Rightarrow$ $d\,o_3=-\frac\epsilon2\{j,\{j,j\}\}-\{j,o_2\}+d\theta$. Exact identities: $\{D,\{D,D\}\}=0$ ($\fsl_N$); $\{j,\{j,j\}\}=0$ for $j=2D+3S_1+5S_0$ ($\fgl_N$); $\{J_T,o'\}=\{J_T,F_{133}\}=\{J_T,F_{223}\}=0$ ($T=d,s_1,s_0$); $d\theta'(T)=0$; $d\{D,\iota\}=0$ $\Rightarrow$ $d\,o_3=0$. The identity is not a test of the kernels: $\{j,o_2\}=0$ identically, every piece $\alt(K^{(v)}_GW^{(v,d)}_G)$ is closed, and the $SL_2$-invariant $\mathsf C$-odd subcomplex at $(3,3)$ has $\dim C^3=1$ ($\Dthree$), $\dim C^4=1$, ${\rm rank}\,d_3=1$ ($\fsl_N$), resp.\ $3,2,2$ ($\fgl_N$): $C^4_{{\rm inv},-}(3,3)=d\,C^3_{{\rm inv},-}(3,3)$.
\end{proposition}
\cert{C44}

\subsection{\texorpdfstring{The class $[O_3]$ for $\fsl_N$}{The class [O3] for sl(N)}}

\begin{theorem}[order $\hbar^3$, $\fsl_N$, $N\ge3$]\label{thm-o3sl}
(i) Weight $(3,3)$, $G$-invariant cochains: $\dim C^2=0$, $\dim C^3=6$, $\dim C^4=39$ ($N\ge4$; $\fsl_3$: $35$), ${\rm rank}\,d_3=6$, ${\rm rank}\,d_4=32$ ($\fsl_3$: $28$); $H^3_{(3,3)}=0$, $H^4_{(3,3)}=\mathbb C[h]$, $h$ double-trace, $\mathsf C$-even; $\dim H^4$ at $(3,3),(4,2),(5,1),(6,0)$: $1,1,0,0$, so $H^4$ of total degree $6$ is the spin-$1$ $SL_2$-module and $(H^4_{(3,3)})^{SL_2}=0$; $\mathsf C$-odd part: $\dim C^3=3$, $\dim C^4=13$, ranks $3,10$, $H^4_-=0$. At $(2,2)$: $H^3=0$, $H^4=1$ ($\fsl_3$: $\dim C^4=14$).
(ii) Primitives:
\begin{equation}
\begin{gathered}\theta'=d\Big(-\frac{N^2}{60}\Dthree\Big),\qquad \{D,\uf\}=d\Big(\frac13\Dthree\Big),\qquad \{D,\iota\}=d\Big(-\frac{3N^2+10}{60}\Dthree\Big),\\ \{D,\iota\}\pm\theta'=d\Big(\mp\frac{N^2}{60}-\frac{3N^2+10}{60}\Big)\Dthree .\end{gathered}
\end{equation}
(iii) $o_3=a\epsilon^2\pi\nu\{D,\iota\}+a\pi\nu\theta'=a\pi\nu(\{D,\iota\}+\theta')=a\pi\nu\,d\xi$, $\xi=-\frac{2N^2+5}{30}\Dthree$; $[O_3]=-\frac{aU}{64\pi^4}([\{D,\iota\}]+[\theta'])=0$, independently of $U$, $\epsilon$, the sign of $a$ and the relative sign of $\theta$ and $\{j,i_2\}$.
(iv) $d\,i_3=\{j,i_2\}+\epsilon\theta=a\epsilon\pi\nu\,d\xi$:
\begin{equation}
i_3=a\epsilon\pi\nu\,\xi=a\,\epsilon U\,\frac{2N^2+5}{1920\pi^4}\,\Dthree,\qquad I_3=a\,\epsilon U\,\frac{2N^2+5}{1920\pi^4}\,\frac13\int dz_1dz_2\,d_{abc}\big(\partial_{z_1}^3A^a\partial_{z_2}^3A^b-3\partial_{z_1}^2\partial_{z_2}A^a\,\partial_{z_1}\partial_{z_2}^2A^b\big)A^c,
\end{equation}
unique ($H^3_{(3,3)}=0$, $C^2_{(3,3)}=0$; in general up to $Q$-exact terms).
\end{theorem}
Symmetry form: $j$ is $\mathsf C$-odd, $i_2$ $\mathsf C$-even, $\theta'$ $\mathsf C$-odd, so $o_3$ is $\mathsf C$-odd; $H^4_{(3,3)}$ is $\mathsf C$-even $\Rightarrow$ $[O_3]=0$, without using $SL_2$; with $SL_2$, also $(H^4_{(3,3)})^{SL_2}=0$. The kernel values enter only through closedness. Coefficients: $N=3,4,5$: $i_2=\epsilon U\frac{37}{1280\pi^4}\uf,\frac{29}{640\pi^4}\uf,\frac{17}{256\pi^4}\uf$; $i_3/a=\epsilon U\frac{23}{1920\pi^4},\frac{37}{1920\pi^4},\frac{11}{384\pi^4}$ times $\Dthree$. The formal identities have coefficients Laurent polynomial in $N$ and hold at each $N\ge3$; $\fsl_3,\fsl_4,\fsl_5$ computed also in faithful coordinates.
\cert{C43, C45}

\subsection{\texorpdfstring{$\fgl_N$, $\kappa=\tr$}{gl(N), kappa = tr}}

\begin{theorem}[order $\hbar^3$, $\fgl_N$, $SL_2$-invariant $I_2$]\label{thm-o3gl}
$j=aD+bS_1+cS_0$. (i) $C^3_{\rm inv}(2,2)=\mathbb C\uf$, $d\uf=\beta\neq0$ $\Rightarrow$ $(H^3_{(2,2)})^{SL_2}=0$; an $SL_2$-invariant solution of the order-$\hbar^2$ equation is
\begin{equation}
i_2=k_2\,\uf,\qquad k_2=\frac{a^2+2ac+b^2+Nbc}2-\frac{\epsilon\pi\nu N^2}{40}\qquad(\text{Theorem~\ref{thm-quad}, }o'=\tfrac1{48}F_{223}-\tfrac{N^2}{40}\beta).
\end{equation}
(ii) $H^4_{\rm inv}(3,3)=1$ ($\dim C^3_{\rm inv}=3$, $\dim C^4_{\rm inv}=11$, ranks $2,8$), quadruple-trace, $\mathsf C$-even. Primitives: $\{D,\uf\}=d(\frac13\Dthree)$, $\{S_1,\uf\}=d(\frac13S_{1,3})$, $\{S_0,\uf\}=0$, $\theta'(d)=d(-\frac{N^2}{60}\Dthree+\frac N{30}S_{1,3})$, $dS_{0,3}=0$.
(iii) $[O_3]=0$ for all $(a,b,c)$, in particular for every $j$ solving the order-$\hbar^2$ equation (Theorem~\ref{thm-Q1}(ii)) and for $s\,j_M+t\,\pi^*D_{\fsl}$;
\begin{equation}
i_3=\frac{k_2}3(a\Dthree+bS_{1,3})+\epsilon\pi\nu\,a\Big(-\frac{N^2}{60}\Dthree+\frac N{30}S_{1,3}\Big)+z\,S_{0,3},\qquad z\text{ free}.
\end{equation}
\end{theorem}
Computed with formal $N$ (trace basis, faithful for $N\ge5$); the identities $O_3=d(\cdot)$ are identities of trace-basis cochains with coefficients polynomial in $N$ and hold at every $N$.
\cert{C46}

\begin{proposition}[$\fgl_N$, non-invariant $I_2$]\label{prop-o3glgen}
Formal $N$ (faithful for $N\ge5$): $H^3_{(2,2)}=Z^3_{(2,2)}=\mathbb Ch_1\oplus\mathbb Ch_2$ ($\dim C^3=8$, ${\rm rank}\,d_3=6$), $h_1$ of type $\tr\cdot\tr$, $h_2$ of type $\tr\cdot\tr\cdot\tr$, both $\mathsf C$-odd and not $SL_2$-invariant. $H^4_{(3,3)}$: $\dim C^3=17$, $\dim C^4=101$, ${\rm rank}\,d_3=13$, ${\rm rank}\,d_4=79$, $\dim H^4_{(3,3)}=9$, $\mathsf C$-even. With classes $e_1,u,v,w\in H^4_{(3,3)}$ (linearly independent):
\begin{gather}
[\{D,h_1\}]=-3e_1,\quad[\{S_1,h_1\}]=-\tfrac{3N}2e_1-2v,\quad[\{S_0,h_1\}]=-Nv,\\ [\{D,h_2\}]=2u,\quad[\{S_1,h_2\}]=Nu+2w,\quad[\{S_0,h_2\}]=Nw;\\
i_2=k_2\uf+y_1h_1+y_2h_2:\quad [O_3]=\epsilon\Big(y_1\big[-\tfrac32(2a+Nb)e_1-(2b+Nc)v\big]+y_2\big[(2a+Nb)u+(2b+Nc)w\big]\Big),
\end{gather}
$=\epsilon(y_1[-\frac32N\chi e_1-(N\gamma-\chi)v]+y_2[N\chi u+(N\gamma-\chi)w])$; $[O_3]=0\iff y=0$ or $\chi=\gamma=0$; for $s\,j_M+t\,\pi^*D_{\fsl}$: $[O_3]=\epsilon s(-\frac32y_1e_1+y_2u)$. $N=3,4$ not computed.
\end{proposition}
\cert{C47}

\subsection{\texorpdfstring{$U(2)$-equivariant quantizations through order $\hbar^3$}{U(2)-equivariant quantizations through order hbar3}}

\begin{theorem}[existence and uniqueness]\label{thm-main}
Setting of Definition~\ref{def-setting}, $U(2)$-equivariant; $U=\pm1$ (Section~\ref{sec-wick}), $\epsilon=U_{\rm tree}$ (Proposition~\ref{prop-tree}), $c_O=-U/(3072\pi^4)$, $\pi\nu=-U/(64\pi^4)$.
\begin{center}
{\small
\begin{tabular}{@{}>{\raggedright\arraybackslash}p{3.0cm}>{\raggedright\arraybackslash}p{1.3cm}>{\raggedright\arraybackslash}p{5.7cm}>{\raggedright\arraybackslash}p{5.4cm}@{}}
\toprule
$\g$ & $[O_2]$ & through $\hbar^2$ & through $\hbar^3$\\
\midrule
$\fsl_N$, $N\ge3$ ($\kappa=\tr$) & $\ne0$ & iff $j=aD$ $+$ exact, $a^2=\epsilon U\frac{N^2}{256\pi^4}$; $i_2$ of Proposition~\ref{prop-I2}, unique & yes, for both signs of $a$; $[O_3]=0$; $i_3$ of Theorem~\ref{thm-o3sl}, unique\\
$\fgl_N$, $N\ge3$, $\kappa=\tr$ & $\ne0$ & iff $j=\alpha\pi^*D_{\fsl}+\chi M_1+\gamma S_0$ $+$ exact, $\chi(N\gamma-\chi)=0$, $\frac{\alpha^2}{12N^2}-\frac{\chi^2}{48}=-\epsilon c_O$; $i_2=k_2\uf$ & yes, for every such $j$; $[O_3]=0$; $i_3$ of Theorem~\ref{thm-o3gl} (free $S_{0,3}$ term)\\
$\fsl_2$, $\fgl_1$, $\fgl_2$ & $0$ & yes; $J_1=0$ admissible (all $J_1$: Theorem~\ref{thm-Q1}(iii)) & Derived statement~\ref{ds-trivial}\\
simple, $\ne\fsl_2$, $\dim(\g^{*\otimes3})^\g=1$ & $\ne0$ & no (any $J_1\in\Vht$, also without equivariance) & ---\\
\bottomrule
\end{tabular}}
\end{center}
\end{theorem}
Proof. Rows 1--2: Theorems~\ref{thm-class-slN}, \ref{thm-Q1}, Proposition~\ref{prop-I2}, Theorems~\ref{thm-o3sl}, \ref{thm-o3gl}, with $U(2)$-equivariance forcing bidegree $(k,k)$ and $SL_2$-invariance, hence $y=0$ in Proposition~\ref{prop-o3glgen}. Row 3: Theorems~\ref{thm-class-slN}, \ref{thm-simple}, \ref{thm-Q1}(iii). Row 4: Theorem~\ref{thm-nogo}. $N=3,4$ in row 2: identities of trace-basis cochains evaluated at $N$.

Scope. Vertices in $\Vht$ (for $J_1\in\Veta\smallsetminus\Vht$: Derived statement~\ref{ds-homotopy}); $\eta$-equivariance at kernel level (Proposition~\ref{prop-eta}); the classes $[O_2]$, $[O_3]$ are computed in $\Cg$ from holomorphic plane waves, their reading in $H^1(V)$ uses Corollary~\ref{cor-O2} and, for $O_3$, Derived statement~\ref{ds-Vht-eta}; $V$ as definition of translation-invariant local functionals; $\kappa=\tr$ for $\fgl_N$ (other invariant forms on $\fgl_N$ not treated); semisimple non-simple $\g$ not treated; $\fso_N,\fsp_N$ with $N>40$: $[O_2]\neq0$ (Theorem~\ref{thm-allsimple}), no-go conditional on $\dim(\g^{*\otimes3})^\g=1$; order $\hbar^4$ not analysed.

\begin{dstatement}[trivial rows and real coefficients]\label{ds-trivial}
(i) $\fsl_2,\fgl_1,\fgl_2$ with $J_1=0$: $\Theta(0)=0$, $\{J_1,I_2\}=0$, so \eqref{eq-o3} is $\{S_{\rm cl},I_3\}=0$ and $I_3=0$ solves; these theories exist through order $\hbar^3$. $J_1\ne0$ for $\fgl_1,\fgl_2$ at order $\hbar^3$: not computed. (ii) With $\epsilon U=+1$ (Derived statement~\ref{ds-epsU}): $a=\pm\frac N{16\pi^2}$, $I_2=\frac{3N^2+10}{1280\pi^4}F_{\uf}$, $I_3=a\frac{2N^2+5}{1920\pi^4}F_{\Dthree}$, $F_w=\frac1{3!}w(A,A,A)$:
\begin{equation}
S=S_{\rm cl}+\hbar\,aF_D+\hbar^2\frac{3N^2+10}{1280\pi^4}F_{\uf}+\hbar^3a\frac{2N^2+5}{1920\pi^4}F_{\Dthree}+O(\hbar^4)\qquad(\fsl_N).
\end{equation}
\end{dstatement}

%======================================================================
\section{Literature}\label{sec-lit}
%======================================================================

Sources read in their PDF text layers unless stated.

\subsection{Wang--Williams 2407.08667 (v3)}
\begin{itemize}\itemsep1pt
\item Sec.~2.2 (8), (9): $\alpha\in\Omega^*(\mathbb R^{d'})\otimes\Omega^{0,*}(\mathbb C^d)\otimes\g[1]$, $S=\int d^dz\,(\frac12\langle\alpha,d\alpha\rangle+\frac16\langle\alpha,[\alpha,\alpha]\rangle)$, $d+d'$ odd; list item (3): $d'=1$, $d=2$, with the non-commutative $\fgl_k$ deformation of 1610.04144.
\item Sec.~2.3: $Q^{\rm GF}=d^*_{\rm dR}+\dbar^*$, $P_\Phi=\frac12Q^{\rm GF}\Phi$ (17), $\dbar^*=-2\sum_i\partial_{z_i}\iota_{\partial/\partial\zb_i}$, $d^*_{\rm dR}=-\sum_i\partial_{x_i}\iota_{\partial/\partial x_i}$ (21), $H=(2^{d+d'}(\pi t)^{d+d'/2})^{-1}e^{-(2\sum|z_i|^2+\sum x_i^2)/4t}d^d\zb\,d^{d'}x$ (22), (25)--(28).
\item Thm~1.1: UV finite to all orders. Thm~1.2: $d'>1$: anomaly free to all orders. Text after Thm~1.2: $d'=1$ theories quantize over $\mathbb C[[\hbar]]/\hbar^2$; topological-holomorphic CS on $\mathbb R\times\mathbb C^2$ is expected to have a two-loop anomaly; no computation.
\item Prop.~3.23: connected decorated Laman graph, $d'\ge1$, $h_1$ odd $\Rightarrow$ $O_{(\Gamma,n)}=0$. Prop.~3.21: non-Laman $\Rightarrow$ $O=0$; Laman: $O$ local, holomorphic derivatives only. (73), (74): $(d+d')|\Gamma'_0|\ge(d+d'-1)|\Gamma'_1|+d+d'+1$, equality for $\Gamma$. Cor.~3.28 excludes items (3), (5). Remark~1.3: GKW's anomaly-freeness uses a gauge in which the topological and holomorphic directions are treated alike. No cohomological description of the obstruction.
\end{itemize}
Overlap: Prop.~3.23 $\leftrightarrow$ $\rho_1=0$; (74) at $d+d'=3$ $\leftrightarrow$ $|I|=3n+1$, $|V|=2n+2$ at $n=2$. Theorems~\ref{thm-assembly}, \ref{thm-class-slN}, \ref{thm-simple}, \ref{thm-allsimple}: the two-loop obstruction of the theory of WW (9) in the heat-kernel gauge of Section~\ref{sec-gauge}.

\subsection{Costello 1610.04144 (v1)}
\begin{itemize}\itemsep1pt
\item Prop.~1.2.1: twisted M-theory on $TN_k\times\mathbb R^3\times\mathbb C^2$ $\simeq$ non-commutative 5d CS on $\mathbb R\times\mathbb C^2$, $GL(k+1)$. Action (Sec.~1.2): $\frac1{2\delta}\int\Tr A\,dA\,dz_1dz_2+\frac1{3\delta}\int\Tr A^3dz_1dz_2+\sum_{n\ge1}\frac{e^n}{3\delta2^nn!}\epsilon^{i_1j_1}..\epsilon^{i_nj_n}\int\Tr A\,\partial_{z_{i_1}}..\partial_{z_{i_n}}A\wedge\partial_{z_{j_1}}..\partial_{z_{j_n}}A\,dz_1dz_2$. Sec.~5: $f*_eg=fg+e\frac12\epsilon^{ij}\partial_if\partial_jg+\dots$; Sec.~9.1: $[z_1,z_2]=e$; the normalization of Proposition~\ref{prop-moyal}.
\item Thm~9.0.2 ($\mathbb R\times X$, $X$ conical symplectic surface, $\fgl_N$ uniformly in $N$, $e^{-k}$ only with $\delta^k$): no obstructions. Thm~9.0.3 ($\mathbb R\times\mathbb C^2$, uniformly in $N$, $\fsl(2)$-invariant, translation invariant, compatible with the $\mathbb C^\times$ action, $e^{-r}\delta^n$ with $r\le n$): no obstructions, two deformations absorbed into $e$, $\delta$. The regime $e^{-r}\delta^n$, $r\le n$, excludes $e\sim\hbar$; no statement at $e=0$. Sec.~9: negative powers of $e$ admitted to cancel potential higher-loop anomalies.
\item Sec.~9.1: local operators $=C^*(\fgl_N[[z_1,z_2]])$; the $GL(N)$-invariant subcomplex has the same cohomology; descent with $\partial/\partial z_i$ acting trivially; the verbal definition of holomorphic translation invariance makes $\partial/\partial\zb_i$ homotopically trivial. App.~B: at $e=0$ the $k$-loop obstruction-deformation complex on $\mathbb R\times X$ is written, $\omega_{\mathbb R\times X}\otimes_{\mathcal D}\mathcal O(\pi^*J\Omega^{-*}_{\rm hol}[t^{-1}][1])[N][-6(k-1)]$, jets pulled back along $\pi:\mathbb R\times X\to X$; for holomorphically translation-invariant quantizations on $\mathbb R\times\mathbb C^2$ the complex $\mathbb C[\eta_3,\eta_5,..]/\mathbb C\otimes\mathbb C[N]\otimes\mathbb C((e))\otimes\mathbb C[dz_1,dz_2][5]$ is the $e\ne0$ complex. Proof of Lemma~B.0.2, first sentence: Section~\ref{sec-h1}, ``Sources''. No mention of a two-loop obstruction of the commutative theory; no statement on an order-$\hbar$ term cancelling an $\hbar^2$ obstruction.
\end{itemize}
Overlap: $j_M=\frac12D$ is the order-$e$ cochain of the Moyal product in Costello's normalization. At finite $N$, $e=0$, $\hbar$-expansion: $q[s\,j_M]=0$ (Theorem~\ref{thm-quad}); $s\,j_M$ alone does not cancel $[O_2]$.

\subsection{Other 5d CS computations}
\begin{itemize}\itemsep1pt
\item Oh--Zhou 2002.07343 \cite{OZ1}: 0- and 1-loop defect diagrams; Sec.~5.3 diagram bounded ((136), (137)), not evaluated. Oh--Zhou 2103.00963 \cite{OZ2}: propagator (38); $P_{32}\wedge P_{24}\wedge P_{43}=0$ (117), $P_{34}\wedge P_{43}=0$ (118); conjecture (119) for single cycles.
\item Ishtiaque--Jeong--Zhou 2408.15712 \cite{IJZ}: gauge (2.11), propagator (2.12), non-commutative $\fgl(1)$; footnote 14 and (C.1) (double edges vanish) remove two-loop defect diagrams. Ashwinkumar 2408.15732 \cite{Ash}: (3.1), (3.2); leading order in $\hbar$. Gaiotto--Oh 1907.06495 \cite{GO}: defect anomalies only.
\item Gwilliam--Rabinovich--Williams 2107.06734 \cite{GRW}: Thm~5.1 one loop; Prop.~6.4: finite one-loop quantization of 5d CS on $\mathbb R\times\mathbb C^2$. Costello 1705.02500 \cite{C17} Sec.~9.3: one loop.
\item INSPIRE citers of 2407.08667 (14 records) and of 1610.04144 (75 records) were scanned: no computation of the order-$\hbar^2$ obstruction of 5d CS on $\mathbb R\times\mathbb C^2$ found.
\end{itemize}
Cycle lengths: $\theta(1,3,3)$: $4,4,6$; $\theta(2,2,3)$: $4,5,5$; no double edge; (117), (118), (119), (C.1), (3.1), (3.2) do not apply.

\subsection{Local cohomology}
Williams 1809.02661 \cite{Wil} Def.~2.5, Prop.~2.8, eq.~(12), Def.~2.28, Lemma~2.29, Cor.~2.30; Gwilliam--Williams 1810.06534 \cite{GW18} Def.~1.14, App.~A (Def.~A.1, Lemmas~A.2, A.7, A.8, Prop.~A.6); Costello 1303.2632 \cite{C13} Sec.~16.1: $\mathbb C\otimes^{\mathbb L}_{\mathbb C[\partial_z,\partial_w]}C^*_{\rm red}(\g[[z,w,e]])$, stated without proof; Saberi--Williams 1910.04120 \cite{SW} Thm~5.5 ($T=0$); Raghavendran--Saberi--Williams 2111.03049 \cite{RSW} Sec.~4 (descent argument). Comparison with Theorems~\ref{thm-A}, \ref{thm-B}: Section~\ref{sec-h1}.

\subsection{\texorpdfstring{The form $\rho_n$}{The form rho\_n}}
\begin{itemize}\itemsep1pt
\item Brown--Hu--Panzer 2406.12734 \cite{BHP} (v2): $\Pf(M)=(2^nn!)^{-1}\sum_{\pi\in S_{2n}}\sgn\pi\,M_{\pi1\pi2}\cdots$ (2.5); $\Pf(\lambda M)=\lambda^n\Pf(M)$ (P2); Def.~2.1 (2.6): $\phi_X=(\det X)^{-1/2}\Pf(dX\,X^{-1}dX)$, $\phi=0$ for $n$ odd; Def.~5.4: $\phi_G=\phi_{\Lambda_C}$, (5.4) sign $\det P$ under change of cycle basis; Lemma~5.5 (5.5) dipole.
\item Balduf--Hu 2503.09558 \cite{BH}: Def.~2: $\phi_G=(-2\pi)^{-l/2}\Pf(d\Lambda_C\Lambda_C^{-1}d\Lambda_C)/\sqrt{\det\Lambda_C}$, $\Lambda_C=C^{\mathsf T}DC$; (1.5): $\phi_{D_3}=\frac1{2\pi}\frac{a_1da_2da_3-a_2da_1da_3+a_3da_1da_2}{(a_1a_2+a_2a_3+a_1a_3)^{3/2}}$; Thm~1: $\alpha_G=\det(C|P)2^{-l}\phi_G$; Lemma~3(ii): $\phi_G=0$ for odd $l$.
\item Balduf--Gaiotto 2408.03192 \cite{BG}: (4.7), Thm~1 (spanning-tree form of $\alpha_\Gamma$), Lemma~4 ($\alpha=0$, odd $L$), Thm~2 ($\alpha\wedge\alpha=0$ for non-trees).
\item GKW 2403.13049 \cite{GKW} (TeX), eq:omegafactorize: $\omega_\Gamma=\alpha_\Gamma^T\prod_{i=1}^H\rho_\Gamma[\lambda^i;z^i]$, $\alpha_\Gamma=\int\prod_e\pi^{-1/2}e^{-s_e^2}ds_e$, $s_e=t_e^{-1/2}(x_{e(0)}-x_{e(1)})$.
\end{itemize}
\begin{dstatement}[$\rho_n$ and $\phi$]
$M=CtC^{\mathsf T}=\Lambda_C$; $\rho_n(M)=(-\frac12)^{n/2}\phi^{\rm BHP}_M=\pi^{n/2}\phi^{\rm BH}_G=\det(C|P)\,(2\sqrt\pi)^n\alpha_G$ for even $n$, up to the cycle-basis sign (BHP (5.4)). At $n=2$, $M=\big(\begin{smallmatrix}x_1+x_3&x_3\\x_3&x_2+x_3\end{smallmatrix}\big)$: $\det M=\Psi$, $\rho_2=-\frac12\phi^{\rm BHP}$, $\rho_2=\pi\phi^{\rm BH}$ (single computation, X1). $n\ge4$: from (P2) only.
\end{dstatement}

\subsection{Odd-loop vanishing}
$T=1$, $n$ odd: reflection (Section~\ref{sec-reg}); WW Prop.~3.23; GKW sec:projective ($x\mapsto-x$ reflects an odd number of $\frac{ds_e}{dt_e}dt_e$, TeX); Balduf--Gaiotto Lemma~4; Balduf--Hu Lemma~3(ii); Budzik--Gaiotto--Kulp--Wu--Yu 2207.14321 \cite{Op}, subsection ``Three-Dimensional Holomorphic-Topological (HT) Twist'' ($\omega^{3d}_\Gamma=\rho\wedge\eta$, odd loops vanish; TeX). Present form: $\rho_{2k+1}=0$ in \eqref{eq-rho}.

\subsection{GKW bitriangle and earlier values}
GKW 2403.13049 (TeX, sec:projective), bitriangle $=\theta(1,2,2)$, edges $t_{12},t_{23},t_{24},t_{31},t_{43}$:
\begin{gather}
\alpha=\frac1{8\pi}\big[t_{31}(t_{24}+t_{43})+t_{12}(t_{23}+t_{24}+t_{43})+t_{23}(t_{24}+t_{31}+t_{43})\big]^{-3/2}\Big[(t_{24}+t_{43})(dt_{12}dt_{23}+dt_{31}dt_{23})\notag\\
+(t_{12}+t_{31})(dt_{23}dt_{24}+dt_{23}dt_{43})-t_{23}(dt_{12}dt_{24}+dt_{12}dt_{43}+dt_{31}dt_{24}+dt_{31}dt_{43})\Big],\qquad\alpha^2=0,\\
H=T=1:\ \ \omega=-\frac{\pi^2}4\,U^{-5/2}\,(t_{12}t_{23}t_{31}t_{24}t_{43})\,e^{-\lambda A^{-1}Bz}\,\lambda_1(\lambda_1+\lambda_2+\lambda_3)\,d\mathrm{Vol}_t,
\end{gather}
no integrated value; GKW Introduction: no non-renormalization theorem for $T=1$; GKW sec:projective: the non-renormalization statement is for $T\ge2$. In the present normalization: $K_\htk(\theta_{122})=-\frac\pi3\lambda_{B_1}\lambda_{C_1}$ at $(H,T)=(1,1)$ (Theorem~\ref{thm-closed}).
\begin{dstatement}[bitriangle]
In GKW's edge list the trivalent vertices are $2,3$; at $\sum\lambda=0$, $\lambda_1(\lambda_1+\lambda_2+\lambda_3)=-\lambda_1\lambda_4$, the product of the bivalent-vertex sources, as $D_\theta=\lambda_{B_1}\lambda_{C_1}$; $(t_{12}t_{23}t_{31}t_{24}t_{43})\,d\mathrm{Vol}_t=\pm\iota_E\prod dt_e$; GKW's $\omega$ and Theorem~\ref{thm-dens} at $(1,1)$ have the same $t$-dependence $U^{-5/2}$; the $t$-independent factor is not compared.
\end{dstatement}
Earlier values ($H=0$): BHP Thm~1.3: $I_{D_3}(1)=1$; Balduf--Hu Sec.~2.2: $\int\phi_{D_{2i+1}}=1$, $I_{D_{2i+1}}=2^{-2i}$; GKW sec:exampleH0T1: $I_N=2^{1-N}$, $N$ odd; here $K_\htk(\theta_{111})=\pi$. Two-loop theta values at $T=0$, $H=2$: S1. No closed value of a two-loop graph integral with $T=1$, $H\ge1$ found in published sources.

%======================================================================
\appendix
\section{Computational certificates}\label{app-cert}
%======================================================================

Supplement folder \path{supplement/}; scripts are named relative to it, \path{supplement/MANIFEST.txt} lists them. ``Exact'': Python integers, \texttt{Fraction}, sympy rationals, or Laurent polynomials in $N$ with rational coefficients. ``Certified rank'': elimination modulo $2147483629$ and $2147483587$, an $r\times r$ minor with nonzero exact (Bareiss) determinant, and exact kernel vectors checked on all rows. ``mod $p$'': computed modulo the stated primes only. Floating point is marked ``float'' with the agreeing digits. Each certificate is an independent recomputation of the statement; where two computations are listed, ``first'' is the original computation and ``recomputed'' the independent one.

{\footnotesize
\renewcommand{\arraystretch}{1.2}
\setlength{\tabcolsep}{3pt}
\begin{longtable}{@{}>{\raggedright\arraybackslash}p{0.6cm}>{\raggedright\arraybackslash}p{2.2cm}>{\raggedright\arraybackslash}p{5.0cm}>{\raggedright\arraybackslash}p{4.4cm}>{\raggedright\arraybackslash}p{3.3cm}@{}}
\toprule
No. & statement & computation, method & exact output & scripts\\
\midrule\endhead
C1 & Prop.~\ref{prop-rules}, engine & recomputed: implementation without engine import; Kirchhoff currents nodal and mesh; $T=1$ Pfaffian from \eqref{eq-rho}; transformation rules & 138 PASS, 0 FAIL; densities 16/16 pointwise exact; rules exact on 11 graph/$(H,T)$ pairs, exact kernels on 9 graphs; $D_\psi=\sigma_T^H\epsilon(n,H)D_{dy}$ on 45 graphs & \path{engine/verify_independent.py}, engine \path{engine/htk/}\\
C2 & Thm.~\ref{thm-dens} & recomputed: KCL$+$KVL with BFS cycle basis, forms evaluated on vectors, exact rationals; 32 graphs ($T=0$: $H=1..6$; $T=1$: $H=0..5$), 3 random rational points, 2 transforms; 53 unsorted $(p,q,r)$ & all equal & \path{theta/vcore.py}, \path{theta/v1_pointwise.py}, \path{theta/v1b_unsorted.py}\\
C3 & Thm.~\ref{thm-intrep} & chart densities exact; $z\ne0$: direct Schwinger integral vs representation (float) & 32/32; 13.2--15.0 digits (5 cases) & \path{theta/v1_pointwise.py}, \path{theta/v4_numerics.py}\\
C4 & Thm.~\ref{thm-closed}, tables & exact $c_1$, $c_0$, signs vs exact engine kernels; basis sign; interior roots; float $z=0$ values; table entries re-evaluated from \eqref{eq-T1}, \eqref{eq-T0} & 16/16, 16/16, 32/32; basis sign on 32 tight graphs and 343 graphs $\theta(p,q,r)$, $p,q,r\le7$, 53 orders; roots 32/32; 13.0--15.9 digits (e.g.\ $0.698131700797724$ vs $2\pi/9$); tables equal & \path{theta/v2_values.py}, \path{theta/v4_numerics.py}, \path{theta/v7_oracle_rerun.py}, \path{arith/arith_check.py}\\
C5 & Laplace expansions & exact, all path orders & 168 graph/$T$ cases, $D\le7$: True & \path{theta/v6_expansion.py}\\
C6 & Prop.~\ref{prop-sourcevals} & exact comparison with S1 (5.6), (7.13) & equal & \path{theta/v3_sources.py}\\
C7 & Prop.~\ref{prop-gauge} & recomputed: sympy exterior algebra, Wirtinger generators; flux by $24^4$ Gauss--Legendre (float) & identities True; flux $1.000000000000008$ at $R=0.7,1.3$; WW (22) factor 4 & \path{normalization/v1_edge_forms.py}\\
C8 & Prop.~\ref{prop-raw} & recomputed: every term of $\bigwedge_e\prod_a(dX^a_e-c_aJ^a_edt_e)$, Gaussian moments, exact at random rational $(t,\lambda)$ & $\Phi_j/(D_\htk(t)(-1)^jt_j)=-1/(64\pi^5)$ at every component $j$; 3 points per theta, 2 per rooted theta & \path{normalization/v2_pointwise.py}, \path{normalization/v2b_rootN.py}\\
C9 & chain of Prop.~\ref{prop-raw} & end-to-end heat-edge position integrals (float) & $(0,1)$ banana(3), banana(5): $0.250000000000000$, $0.0625$; $(1,0)$ banana(2), (3): 15.6--17.0 digits & \path{normalization/v4_calibration.py}\\
C10 & Prop.~\ref{prop-wick} & own permutation signs; $\sigma C_{L'}$ vs index forms for $\fso(3)$, $\fso(4)$; 5 random relabellings per graph with controls & $\sigma=-1,-1$; $\sigma C_{L'}=W^\Gamma$; invariance 5/5; controls fail & \path{normalization/v3_wick.py}, \path{anomaly/v_orientation.py}\\
C11 & Thm.~\ref{thm-assembly}, Prop.~\ref{prop-count} & sympy assembly; symbolic count; tight-graph catalog & $O_2/(\Delta(-3W_{133}+2W_{223}))=-U\hbar^2/(15360\pi^4)$; $[O_2]=-U\hbar^2/(3072\pi^4)[F_{223}]$; count $n=1..6$; $D=3$: two graphs, $E=4$; $|\Aut|=4,4$ & \path{normalization/v5_assembly_count.py}, \path{anomaly/v_b1.py}\\
C12 & Thm.~\ref{thm-kernels} & first: exact chart identity at 5 points with sources at all non-root vertices, float integrals; recomputed: own Kirchhoff and forms from \eqref{eq-K}, exact density, mpmath tanh-sinh (dps 30) & $U=1/(w_Aw_Bw_C)$; $D\,U^{7/2}/([p_1p_2][p_3p_4])=-\frac12,+\frac12$ (8/8), independent of $\lambda_N$; integrals: relative differences $2.1\cdot10^{-17}$, $8.2\cdot10^{-18}$ from $-\pi/20$, $\pi/30$ & \path{anomaly/v_kernels.py}, \path{referee/kern_check.py}\\
C13 & Thm.~\ref{thm-colour} & index forms contracted with exact $\kappa^{-1}$ vs trace formulas; own Casimir contraction & $\fgl_{2,3,4}$, $\fsl_{2,3,4}$: True; $\fsl_3$ at 10 points: True; $\fsl_2$ form True & \path{anomaly/v_lie_ce.py}, \path{deform/t_weights.py}, \path{referee/ce_check.py}\\
C14 & Prop.~\ref{prop-closed} & formal $N$: exact; fixed algebras: evaluation at random momenta and coordinates (8 algebras, range 30), 24 five-momentum configurations (12 constructions), random points mod $p$ up to $E_8$; controls & all zero; controls $d\beta=0$, $d(\text{random})\ne0$, permuted coefficients $d\neq0$ & \path{deform/vformal.py}, \path{anomaly/v_ce_large.py}, \path{simpleg/t_closed.py}, \path{simpleall/va_points.py}, \path{simpleall/t_dF_control.py}\\
C15 & Thm.~\ref{thm-class-slN} & first: sample ranks; recomputed: exact elimination over $\mathbb Q$, fixed-$N$ faithfulness by exact Gram ranks ($\fgl_N$) and binomial inversion ($\fsl_N$); third computation: $\fsl_3$ exact ranks at 50 sample points & formal $36/26/10/6/4$ ($\fgl_N$), $15/11/4/3/1$ ($\fsl_N$); fixed-$N$ table as stated; $\fsl_3$: ${\rm rank}\,d=3\to{\rm rank}[d|F_{223}]=4$, $F_{133}+F_{223}=18\beta_{\tr}$ & \path{anomaly/v_formalN.py}, \path{deform/vformal.py}, \path{deform/vfixed.py}, \path{referee/ce_check.py}\\
C16 & Lemma~\ref{lem-red}, Thm.~\ref{thm-simple} ($\fsl_{3,4,5}$, $\fso_{5,6,7,8}$, $\fsp_{4,6}$, $G_2$, $F_4$) & own constructions ($\fso_N$, $\fsp_{2m}$, $\fsl_N$ matrices, $G_2=\mathrm{Der}(\mathbb O)$, $F_4=\fso_9\oplus S_{16}$); full coefficient 4-tensors; exact membership on all entries or printed $2\times2$ minor; unreduced complex for 4 algebras & minors for $X=F_{223}$: $\fsl_3$ $-1451329941600$, $\fso_5$ $-1926955008$, $G_2$ $-116095057920$, $F_4$ $-19670999040$; $c_{\rm Killing}=\frac12$ on all entries (43869696 for $F_4$); rescaling $c=2,3$ exact & \path{simpleg/t_classes.py}, \path{simpleg/t_classes_big.py}, \path{simpleg/t_nonreduced.py}\\
C17 & Prop.~\ref{prop-inv} (table) & own Weyl formula; nullity modulo $1000003$, $999983$ on own constructions & tables identical & \path{simpleg/t_inv_weyl.py}, \path{simpleg/t_inv_lin.py}\\
C18 & Thm.~\ref{thm-simple} (all of $\mathcal L$) & first: exact, all entries of the full coefficient tensors (Kac constructions; $E_8$: $1.6\cdot10^6$ entries); recomputed: own constructions ($G_2$ in $\fso_7$, $F_4$ via Clifford module, $E_{6,7,8}$ Frenkel--Kac), Jacobi and invariance exact; random points mod $1048573$, $1048571$ (23 algebras); all entries mod $524287$, $524269$ (21 algebras; $\fso_{12}$, $E_6$ one prime); & $F_{133}+F_{223}=\frac12\beta_K$ in all cases; 220 nonzero minors for $\g\ne\fsl_2$ (e.g.\ $E_8$ $F_{223}$ minor $924472$ mod $1048573$; first computation $46080$); $\fsl_2$: $F_{223}=-3\beta$ & \path{simpleall/run_a.py}, \path{simpleall/va_points.py}, \path{simpleall/vb_full.py}, \path{simpleall/vlie.py}\\
C19 & Prop.~\ref{prop-inv} (further) & first: nullity mod $1000003$, $999983$ of invariance systems, explicit invariants as lower bounds ($E_8$: $S^3$ 3320 unknowns, 18176 equations); recomputed: root data only, Racah--Speiser/Brauer--Klimyk, $\dim(\g^{*\otimes3})^\g={\rm mult}(\g,\g\otimes\g)$ & non-A: ${\rm mult}(\g,\g\otimes\g)=1$, ${\rm mult}(\g,S^2\g)=0$, ${\rm mult}(\g,\Lambda^2\g)=1$ ($\fso_N$, $N\le40$, $N\ne6$; $\fsp_N$, $N\le40$; exceptional); $\fsl_3..\fsl_{30}$: $2,1,1$ & \path{simpleall/inv3.py}, \path{simpleall/vbk.py}\\
C20 & $F_{223}=x\beta+aG_{kk}$, $dG_{kk}=0$ & first: exact, all entries, 17 algebras; recomputed: random points mod two primes (22 algebras), all entries (19); $dG_{kk}$ symbolic (sympy) & identity with the tabulated $(x,a)$; $G_{kk}\notin\mathbb C\beta$; the 10 polynomials $P_{lm}$ of $dG_{kk}=\sum P_{lm}\kappa_{lm}f_{abc}$ vanish identically & \path{simpleall/decomp_check.py}, \path{simpleall/dGkk.py}, \path{simpleall/va_points.py}, \path{simpleall/v_misc.py}\\
C21 & probe $\varpi$, transfer, networks & first: exact transfer coefficients, oriented-graph isomorphisms (all isomorphisms), networks on 22 algebras; recomputed: dense mod-$p$ tensors on 20 algebras; random totally antisymmetric $f$ with random metric (3 seeds, 2 primes) & $\ell(F_{223})=(-6d,-8n_4)$, $\ell(\beta)=(2d,d)$, $\ell(G_{kk})=(8d(d-1),8d)$, $\varpi(F_{223})=2d(8n_4-3d)$, $\varpi(F_{133})=-\varpi(F_{223})$, $\varpi(G_{kk})=8d^2(d-3)$; network identities hold for random $f$ & \path{simpleall/part_b.py}, \path{simpleall/analysis_b.py}, \path{simpleall/networks.py}, \path{simpleall/vb_dense.py}\\
C22 & Thm.~\ref{thm-allsimple}(i) & first: $n_4$ on matrix families, interpolation through 9 values ($\fso$: $N=5..22$, $\fsp$: $4..24$, $\fsl$: $2..16$); recomputed: Casimir eigenvalues and Weyl dimensions, $n_4=\sum_Vc_V\dim V\,w_V^4$ & polynomials equal at all computed $N$; recomputed range $\fso_N$ $5..40$, $\fsp_N$ $4..40$, $\fsl_N$ $2..30$; exceptional $n_4$: $G_2$ $\frac{119}{96}$, $F_4$ $-\frac{13}{81}$, $E_6$ $-\frac{39}{32}$, $E_7$ $-\frac{2261}{648}$, $E_8$ $-\frac{124}{15}$ & \path{simpleall/family.py}, \path{simpleall/interp.py}, \path{simpleall/vbk_family.py}\\
C23 & Thm.~\ref{thm-allsimple}(iv) & sympy identities; MSV parameters vs computed $d$, $n_4$ & True; 64 algebras & \path{simpleall/vogel_check.py}, \path{simpleall/v_misc.py}\\
C24 & Prop.~\ref{prop-H3all} & first: exact ranks with left kernels (10 algebras up to $E_6$); recomputed: universal computation over $\mathbb Q$ (sympy), numeric mod $p$ at 6 random points on 12 algebras incl.\ $E_7$, $E_8$ & table as stated; $(T_1,T_3)$ value rank 2 for all & \path{simpleall/part_c.py}, \path{simpleall/vc.py}\\
C25 & $dE_f\neq0$ & first: all entries; recomputed: random points mod $p$, compared with the $T_1/T_3$ formula & nonzero entries $\fso_8$ 47088, $G_2$ 8784, $F_4$ 226944; recomputed values $\fso_8$ 494130, $G_2$ 739788, $F_4$ 189015, equal to the formula & \path{simpleall/part_c.py}, \path{simpleall/vc.py}\\
C26 & Thm.~\ref{thm-A} & first: jet model, full differential, complete weight slices, ranks mod $2147483629$, $2147483587$ and exact over $\mathbb Q$; recomputed: own jet model, ranks mod $10^9+7$, $998244353$, exact over $\mathbb Q$ for ab1, ab2, aff1, heis, unrestricted $\fsl_2$; representatives $\psi_S$ of exact $Z^n(\Cg)$ bases & first: 1087 entries, 0 mismatches; recomputed: 2509 distinct entries, 0 mismatches; 376/376 representative tests; largest: $\fgl_2$, $w=0$, $(2,2)$, $g=1$, $\dim M=28566$, $H=3$ (predicted $3$) & \path{h1/check_h1.py}, \path{h1/verify/vjet.py}, \path{h1/verify/vcohom.py}, \path{h1/verify/run_checks.py}\\
C27 & Thm.~\ref{thm-B} & strict invariants $\Veta_t=\{x:\eta_vx\in I_{t+e_v}\}/I_t$ as kernels mod $p$; $\eta_S\phi$ vs $\psi_S(\phi)$ exact & $H^g(\Veta)=H^{g+3}(\Cg)$ at $w=0$, $0$ elsewhere; $\dim\Veta=0$ in 104 negative-weight entries; (B4) for ab1, $\fsl_2$, $\fgl_2$, aff1: 0 mismatches; $\eta_S\phi=\pm\psi_S(\phi)$ & \path{h1/verify/vcohom.py}, \path{h1/verify/eta_psi.py}, \path{h1/check_h1.py}\\
C28 & Thm.~\ref{thm-C} & unrestricted $\fsl_2$ exact over $\mathbb Q$ vs charge-0; $\Cg$ side vs $G$-invariant trace-basis code & equal (104 entries); $\fgl_2$: $H^3_{(1,1)}=2$, $H^4_{(2,2)}=3$; $\fsl_2$: $H^4_{(2,2)}=0$ & \path{h1/verify/cetable.py}, \path{h1/qtable.py}\\
C29 & Prop.~\ref{prop-HCEfalse} & components of $G_\epsilon$, $T_\lambda$ in the superfield algebra, exact over $\mathbb Q$ & closed, rank 1 modulo $B^0$ (ab1, $\fsl_2$, $\fgl_2$; $T_\lambda$: $\fgl_2$) & \path{h1/verify/extra.py}\\
C30 & Prop.~\ref{prop-chars} & enumeration of $S$, $k=0..5$, $0\le a,b\le2k+3$ & only $(k,S=U,a=b=k)$ & \path{h1/verify/misc.py}\\
C31 & Prop.~\ref{prop-eta} & sympy operator identities; $(\iota^x+\iota^y)$ on $k(x-y)\prod(dx-dy)$ and on $D^*_x$ of it, control $\iota^x-\iota^y$; $\eta_v(6I)$ in $M$ exact & identities True; control nonzero; $\eta_v(6I)=0$ & \path{h1/verify/misc.py}, \path{h1/verify/extra.py}\\
C32 & Cor.~\ref{cor-O2}, Prop.~\ref{prop-J1eta} & $\fsl_2$, $\fgl_2$: $S_{\rm free}$, $I$ modulo $B^0$; $\eta_vI$, $\eta_vS_{\rm free}$ modulo $DM$; $\fgl_2$ slice $w=0$, $(1,1)$, $g=0$ & $2S_{\rm free}+3I\in B^0$; $\eta_vI\in DM$, $\eta_vS_{\rm free}\notin DM$; $\dim M=1102$, $H=2$ (predicted 2), representatives pass & \path{h1/verify/extra.py}, \path{h1/verify/run_checks.py}\\
C33 & \eqref{eq-jetid} & 150 monomials per algebra (ab2, $\fsl_2$, aff1, heis, $\fgl_2$); control: one Koszul sign removed & all True; control $Q^2\ne0$ on 271 of 418 monomials & \path{h1/verify/test_identities.py}\\
C34 & Props.~\ref{prop-qme}, \ref{prop-o3eq} & engine kernels with current insertions; brute-force top forms; controls & polygon(4) at $(2,1)$: $K=0$ for 5 insertion patterns (control at $(3,0)$: $\frac{29}{36}$); 16 tight three-loop graphs: density $0$ at 2 random points each (theta controls nonzero) & \path{o3/verify/v7b_controls.py}, \path{o3/verify/v7c_threeloop.py}, \path{o3/o0_vanishing.py}\\
C35 & Prop.~\ref{prop-cochain} & own engine; formal vs explicit dual-basis on $\fgl_2,\fgl_3,\fsl_2,\fsl_3$ & all identities True & \path{deform/t_engine.py}, \path{o3/verify/v1_libtests.py}\\
C36 & Thm.~\ref{thm-H3} & formal and fixed-$N$ ranks; $\fgl_2$ relation symbolic; $dJ_1=0$ analytic; random symmetric $T$ control & dimensions as stated; $dJ\ne0$ for random $T$ & \path{deform/vformal.py}, \path{deform/vfixed.py}, \path{deform/t_explicit.py}, \path{simpleg/t_j1.py}, \path{simpleg/t_sl2inv.py}\\
C37 & Prop.~\ref{prop-moyal} & formal; explicit $\fgl_2$, $\fgl_3$ at 3 random points & True; $\{j_M,j_M\}\neq0$ & \path{deform/t_explicit.py}, \path{o3/verify/v1_libtests.py}\\
C38 & Thms.~\ref{thm-quad}, \ref{thm-Q1} & formal Laurent identities in $N$; fixed $N=3,4,5$; explicit $\fsl_3$ coboundary; exact witnesses; coefficient arithmetic & True; $\pi^*D_{\fsl}$: $(F_1{:}\,0,\ F_{223}{:}\,\frac1{108},\frac1{192},\frac1{300})$; $S_1-\frac2NS_0$: $(0,-\frac1{48})$; $j_M$: $(0,0)$; explicit $\fsl_3$: $\frac12\{D,D\}=F_{223}/108+d[\frac{13}{18}\uf]$ at 50 points & \path{deform/vformal.py}, \path{deform/vfixed.py}, \path{deform/t_coeff.py}, \path{referee/ce_check.py}\\
C39 & Props.~\ref{prop-tree}, \ref{prop-o3rules}(i) & exact; explicit $\fgl_3$ dual basis, random integer points, all 36 pairs of local orders & $\alt(\frac13p_1^1p_2^2d)=D$ ($\fgl$, $\fsl$); tree identity True for the four pairs (e.g.\ $1152=1152$) & \path{o3/verify/v7_misc.py}, \path{o3/o10_tree_bracket.py}\\
C40 & Prop.~\ref{prop-o3rules}(ii) & $H=1$ position integrals on bananas (scipy dblquad, float) vs $\mp y\lambda$ & recomputed: 8 comparisons, 15.3--16.3 digits; first: 24 comparisons, 13.5--15.9 digits; opposite-sign control: relative difference $2.00$ & \path{o3/verify/v9_extra.py}, \path{o3/o1_banana_derivative.py}\\
C41 & Prop.~\ref{prop-o3kernels} & spred reduction with own insertion, brute-force density checks; oracle (lattice rule, no reduction code, float) & all 12 rational multiples of $\pi$, equal up to the uniform local-order sign; $\theta(1,3,3)$ $v=1$: $-\frac{101\pi}{180}=-1.7627825$, oracle $-1.7625165\pm2.0\cdot10^{-3}$; $\theta(2,2,3)$ $v=5$: $\frac\pi{15}=0.2094395$, oracle $0.2094172\pm5.7\cdot10^{-4}$; $\theta(2,2,3)$ $v=3$: $-\frac{8\pi}{27}$, oracle $-0.9307405\pm1.1\cdot10^{-3}$; $\theta(1,3,3)$ $v=0$: $\frac{35\pi}{324}$, oracle $0.3393379\pm8.6\cdot10^{-4}$ & \path{o3/verify/v2_kernels.py}, \path{o3/verify/v8_oracle.py}, \path{o3/o2_kernels.py}, \path{o3/o16_oracle.py}\\
C42 & Prop.~\ref{prop-theta} (colour) & double-line contraction (formal $N$) vs explicit einsum with exact $\kappa^{-1}$ & recomputed: $\fgl_3$ 16/16, $\fsl_3$ 10/10; first: $\fgl_3,\fsl_3,\fgl_2$ 36/36 & \path{o3/verify/v3_assemble.py}, \path{o3/o18_colour.py}, \path{o3/o19_colour_check.py}\\
C43 & Props.~\ref{prop-I2}, \ref{prop-theta}, Thm.~\ref{thm-o3sl}(ii)(iii) & exact formal-$N$ identities; $d\Dthree$ formula at 4 points on $\fgl_3$, $\fsl_3$; ranks at $N=3,4,5,7,11$ & $o'-\frac{N^2}8\{D,D\}=d\iota$; $\theta'(d)$ as stated; $\theta'+\{D,\iota\}=d(-\frac{2N^2+5}{30}\Dthree)$; $\theta'-\{D,\iota\}=d(\frac{N^2+5}{30}\Dthree)$; $C^2_{(2,2)}=0$, $\dim C^3_{(2,2)}=3$, ${\rm rank}\,d_3=3$ & \path{o3/verify/v3_assemble.py}, \path{o3/verify/v7_misc.py}, \path{o3/o4_i2_sl.py}, \path{o3/o3_theta.py}, \path{o3/o14_pieces.py}, \path{o3/o17_final.py}\\
C44 & Prop.~\ref{prop-o3closed} & exact cochain identities; subcomplex ranks & all identities True; controls $\{D,r\}$: 10320 ($\fsl$), 12480 ($\fgl$) nonzero coordinates; $\mathsf C$-odd $SL_2$-invariant $(3,3)$: $1,1,1$ ($\fsl_N$, $N=3,4,5,7,11$), $3,2,2$ ($\fgl_7$) & \path{o3/verify/v3_assemble.py}, \path{o3/verify/v5_cohom_sl.py}, \path{o3/verify/v9_extra.py}, \path{o3/o21_brackets_F.py}\\
C45 & Thm.~\ref{thm-o3sl}(i) & first: certified ranks (minor determinants 324, 63700992), formal $N$; recomputed: exact over $\mathbb Q$ at $N=7,11$ (trace basis, faithful) and $N=3,4,5$ (faithful coordinates); classes by ${\rm rank}[B^4|c]$ vs ${\rm rank}\,B^4$ & as stated; $\fsl_3$: $(2,2)$: $3,14/3,10$; $(3,3)$: $6,35/6,28$; $\theta'$, $\{D,\iota\}$, $o_3'$ exact at $N=3,4,5$; $h-\mathsf Ch\in B^4$ & \path{o3/verify/v4_faithful.py}, \path{o3/verify/v5_cohom_sl.py}, \path{o3/o5_class_sl.py}, \path{o3/o7_sl2_structure.py}, \path{o3/o8_parity.py}, \path{o3/o9_fixedN.py}\\
C46 & Thms.~\ref{thm-quad}, \ref{thm-o3gl} & exact formal-$N$ decompositions in $\{F_0,F_1,\beta\}$; $\fgl$ cohomology at $N=5,7,11$ & decompositions True; $H^4_{\rm inv}(3,3)=1$ ($3,11$; ranks $2,8$); primitives True & \path{o3/verify/v3_assemble.py}, \path{o3/verify/v6_gl.py}, \path{o3/o11_gl.py}, \path{o3/o13_structure.py}, \path{o3/o20_s0.py}\\
C47 & Prop.~\ref{prop-o3glgen} & first: certified ranks (minor determinants 559872, 21035720123168587776); recomputed: exact ranks at $N=5,7,11$ & $H^3_{(2,2)}=2$; $H^4_{(3,3)}=101-79-13=9$, $\mathsf C$-odd part $0$; $[\{T,h_i\}]$ span 4 dimensions; ${\rm rank}(h\mapsto[\{J,h\}])=0$ for $J=\pi^*D_{\fsl}$, $2$ for $j_M,D,S_1,S_0$, $j_M+3\pi^*D_{\fsl}$, random & \path{o3/verify/v6_gl.py}, \path{o3/o12_gl_full.py}\\
\bottomrule
\end{longtable}}

Single computations (no independent recomputation; used only in Derived statements): X1, $\rho_2=-\frac12\phi^{\rm BHP}=\pi\phi^{\rm BH}$ at $n=2$ (exact; \path{lit/rho2_dipole.py}); X2, Derived statement~\ref{ds-C4inv} (\path{simpleg/t_c4inv.py}, \path{simpleg/t_c4span.py}).

Not done: no end-to-end position-space integral of a $(2,1)$ theta, with or without a $J_1$ vertex; the signs $U$ and $s_h$, $s_P$ are not fixed; $\epsilon U=+1$ rests on Derived statement~\ref{ds-epsU}; no Feynman-graph computation of $\eta_vO[L]$; no computation for the Moyal ($e\ne0$) complex; $\fgl_N$ classes with non-invariant $I_2$ at $N=3,4$; order $\hbar^4$.

%======================================================================
\begin{thebibliography}{99}\small
\bibitem{S1} Li, Liu, Chang, \emph{Holomorphic Feynman diagrams in four dimensions and the quantum BPS ring}, arXiv:2611.24160.
\bibitem{Op} K.~Budzik, D.~Gaiotto, J.~Kulp, J.~Wu, M.~Yu, \emph{Feynman Diagrams in Four-Dimensional Holomorphic Theories and the Operatope}, arXiv:2207.14321.
\bibitem{GKW} D.~Gaiotto, J.~Kulp, J.~Wu, \emph{Higher Operations in Perturbation Theory}, arXiv:2403.13049.
\bibitem{WW} Wang, Williams, arXiv:2407.08667 (v3, 13 April 2026).
\bibitem{C16} Costello, arXiv:1610.04144.
\bibitem{C17} Costello, arXiv:1705.02500.
\bibitem{C13} Costello, arXiv:1303.2632.
\bibitem{Wil} Williams, arXiv:1809.02661 (v2).
\bibitem{GW18} O.~Gwilliam, B.~R.~Williams, arXiv:1810.06534 (v2).
\bibitem{SW} Saberi, Williams, arXiv:1910.04120.
\bibitem{RSW} Raghavendran, Saberi, Williams, arXiv:2111.03049.
\bibitem{GRW} Gwilliam, Rabinovich, Williams, arXiv:2107.06734.
\bibitem{OZ1} Oh, Zhou, arXiv:2002.07343.
\bibitem{OZ2} Oh, Zhou, arXiv:2103.00963.
\bibitem{IJZ} Ishtiaque, Jeong, Zhou, arXiv:2408.15712.
\bibitem{Ash} Ashwinkumar, arXiv:2408.15732.
\bibitem{GO} Gaiotto, Oh, arXiv:1907.06495.
\bibitem{BHP} Brown, Hu, Panzer, arXiv:2406.12734.
\bibitem{BH} Balduf, Hu, \emph{The Topological form is the Pfaffian form}, arXiv:2503.09558.
\bibitem{BG} Balduf, Gaiotto, arXiv:2408.03192.
\bibitem{MSV} R.~L.~Mkrtchyan, A.~N.~Sergeev, A.~P.~Veselov, arXiv:1105.0115 (v3).
\bibitem{Kac} V.~G.~Kac, \emph{Infinite-dimensional Lie algebras}, Sec.~7.8 (construction of the simple algebras in C18, C19).
\end{thebibliography}

\end{document}
